\makeatletter
%% ===== 「名前と操作」の型（tools/kata9-int.md、組 V9-2）。頭文字 kixf は章ごとに置きかえる =====
\colorlet{kixfNc}{bkOrange!85!black}% 名前の枠の線（点線）
\colorlet{kixfSc}{bkTheme}% 実際の操作の枠の線（実線）
\colorlet{kixfSb}{bkTheme!8}% 実際の操作の地
\tikzset{
  kixf dan/.style={font=\sffamily\bfseries\footnotesize, rounded corners=1.2mm, inner xsep=1.6mm, inner ysep=0.5mm, line width=0.9pt},
  kixf dan1/.style={kixf dan, draw=kixfNc, dash pattern=on 2.2pt off 1.4pt, fill=white, text=bkInk},
  kixf dan2/.style={kixf dan, draw=kixfSc, fill=kixfSb, text=bkInk},
  kixf dan3/.style={kixf dan, draw=bkTheme, fill=bkTheme, text=bkOn},
  kixf name/.style={draw=kixfNc, line width=0.9pt, dash pattern=on 2.2pt off 1.4pt, rounded corners=1.2mm, inner xsep=1.2mm, inner ysep=1.2mm, fill=white},
  kixf seki/.style={draw=kixfSc, line width=0.9pt, rounded corners=1.2mm, inner xsep=1.2mm, inner ysep=1.2mm, fill=kixfSb},
  kixf tag/.style={font=\sffamily\bfseries\footnotesize, text=bkInk, inner sep=0.4mm},
  kixf q/.style={diamond, aspect=4.6, draw=bkYellow!75!black, line width=0.9pt, fill=bkYellowBg, font=\sffamily\footnotesize, align=flush center, inner sep=0.2mm, text=bkInk},
  kixf do/.style={rectangle, draw=\bkTT, line width=0.9pt, fill=white, rounded corners=0.6mm, font=\sffamily\footnotesize, align=flush center, inner xsep=1.5mm, inner ysep=1.2mm, text=bkInk},
  kixf go/.style={-{Stealth[length=2.2mm]}, line width=0.9pt, bkInk!75, shorten <=1.5mm, shorten >=1.5mm},
  kixf yn/.style={font=\sffamily\bfseries\footnotesize, text=bkInk, fill=white, inner sep=0.5mm},
}
% 実際の操作の4種類の名前（実線の札の字）  1 ふつうの定積分　2 表を引く　3 部分分数　4 行列の計算
\def\kixf@op#1{\ifcase#1\or ふつうの定積分\or 表を引く\or 部分分数\or 行列の計算\fi}
% 段の札（例題の中で）  \kixfdan{1|2|3}
\newcommand\kixfdan[1]{\tikz[baseline=(d.base)]\node[kixf dan#1] (d) {\vphantom{あ}\ifcase#1\or 名前\or 直す\or 楽をする\fi};}
% 直す段の札に、どの道具かをそえる（例題の中で）  \kixfnaosu{1|2|3|4} →「直す：表を引く」
\newcommand\kixfnaosu[1]{\tikz[baseline=(d.base)]\node[kixf dan2] (d) {\vphantom{あ}直す：\kixf@op{#1}};}
% 実線の札だけ（表の「手順」の頭・本文）  \kixfop{1|2|3|4}
\newcommand\kixfop[1]{\tikz[baseline=(d.base)]\node[kixf dan2] (d) {\vphantom{あ}\kixf@op{#1}};}
% 段の見出し（本文で）  \kixfdanhead{1|2|3}{見出し}
\newcommand\kixfdanhead[2]{\par\addvspace{8pt}\Needspace{7\baselineskip}\noindent
  \tikz[baseline=(d.base)]\node[kixf dan#1, font=\sffamily\bfseries\small, inner ysep=0.8mm] (d) {\vphantom{あ}\ifcase#1\or ①\,名前の段\or ②\,直す段\or ③\,楽をする段\fi};%
  \hspace{0.7\zw}{\sffamily\bfseries\color{\bkTT}#2}\par\nobreak\vspace{3pt}\noindent\ignorespaces}
% 場合の札  \kixfcase{C}（′ は \kixfcase{A$'$}）
\newcommand\kixfcase[1]{\tikz[baseline=(c.base)]\node[fill=bkTheme, text=bkOn, font=\sffamily\bfseries\small, rounded corners=0.8mm, inner xsep=1mm, inner ysep=0.6mm, minimum width=5.4mm] (c) {\vphantom{Ag}#1};}
% 例題の1行め  \kixfbaai{ことば}{9-2}{C}（場合が2つなら {D,B}：通った順）
\newcommand\kixfbaai[3]{{\sffamily\bfseries 場合：}#1（表#2の\kixf@cases#3,\@nil）}
\def\kixf@cases#1,#2\@nil{\kixfcase{#1}\ifx\relax#2\relax\else と\kixf@cases#2\@nil\fi}
% 式の中の「名前」（点線）と「実際の操作」（実線）
\newcommand\kixfN[1]{\tikz[baseline=(n.base)]\node[kixf name] (n) {$\displaystyle #1$};}
\newcommand\kixfNt[1]{\tikz[baseline=(n.base)]{\node[kixf name] (n) {$\displaystyle #1$};\node[kixf tag, anchor=north] at ([yshift=-0.8mm]n.south) {名前};}}
\newcommand\kixfS[1]{\tikz[baseline=(n.base)]\node[kixf seki] (n) {$\displaystyle #1$};}
% 札つきの実線の枠  \kixfSt{1|2|3|4}{式}（札の字：ふつうの定積分・表を引く・部分分数・行列の計算）
\newcommand\kixfSt[2]{\tikz[baseline=(n.base)]{\node[kixf seki] (n) {$\displaystyle #2$};\node[kixf tag, anchor=north] at ([yshift=-0.8mm]n.south) {\kixf@op{#1}};}}
% 場合分けの表の題の帯  \kixfcasetitle{9-2}{題}
\newcommand\kixfcasetitle[2]{\par\addvspace{6pt}\Needspace{8\baselineskip}\noindent\begin{tikzpicture}[baseline=(t.base)]
  \node[fill=bkTheme, text=bkOn, font=\sffamily\bfseries\small, rounded corners=1mm, inner xsep=2mm, inner ysep=1mm] (t) {場合分けの表\,#1};
  \node[anchor=west, font=\sffamily\bfseries, text=\bkTT, inner sep=0pt] at ([xshift=2mm]t.east) {#2};
  \end{tikzpicture}\par\nobreak\vspace{3pt}}
% 名前と操作の対応表の小さな絵  \kixfpic{1〜7}
%  1 フーリエ係数（1周期で cos をかける）　2 複素係数（回る矢印）　3 フーリエ変換（−∞〜∞）
%  4 ラプラス変換（0から、e^{-st} で小さく）　5 逆ラプラス変換（s の平面の極）　6 伝達関数（箱）　7 e^{At}（行列とうずまき）
\newcommand\kixfpic[1]{\begin{tikzpicture}[x=1mm, y=1mm, baseline=(current bounding box.center)]
  \ifcase#1\or
  \draw[bkGray, line width=0.4pt] (0,4) -- (14,4);
  \draw[bkInk, line width=0.8pt] (1,4) -- (1,7) -- (7,7) -- (7,1) -- (13,1) -- (13,4);
  \draw[bkTheme, line width=0.8pt, dash pattern=on 1.6pt off 0.9pt, domain=1:13, samples=40, smooth] plot (\x,{4+2.6*cos(30*(\x-1))});
  \draw[bkInk, line width=0.5pt] (1,-0.6) -- (1,-1.6) -- (13,-1.6) -- (13,-0.6);
  \or
  \draw[bkGray, line width=0.4pt] (2,4) -- (12,4) (7,-0.5) -- (7,8.5);
  \draw[bkInk, line width=0.7pt] (7,4) circle (3.6);
  \draw[bkTheme, line width=1.1pt, -{Stealth[length=1.6mm]}] (7,4) -- ({7+3.6*cos(-40)},{4+3.6*sin(-40)});
  \draw[bkInk, line width=0.6pt, -{Stealth[length=1.3mm]}] ({7+5*cos(10)},{4+5*sin(10)}) arc (10:-60:5);
  \or
  \draw[bkInk, line width=0.6pt, {Stealth[length=1.4mm]}-{Stealth[length=1.4mm]}] (0,1) -- (14,1);
  \draw[bkInk, line width=0.8pt] (1.5,1) -- (5.5,1) -- (5.5,6.5) -- (8.5,6.5) -- (8.5,1);
  \draw[bkTheme, line width=0.8pt, dash pattern=on 1.6pt off 0.9pt, domain=1.5:12.5, samples=50, smooth] plot (\x,{1+2.2*(1+cos(80*(\x-7)))/2});
  \or
  \draw[bkInk, line width=0.6pt, -{Stealth[length=1.4mm]}] (1,1) -- (14,1);
  \draw[bkInk, line width=0.6pt] (1,0) -- (1,8.5);
  \fill[bkTheme!25, domain=1:13, samples=40, smooth] (1,1) -- plot (\x,{1+6.5*exp(-0.3*(\x-1))}) -- (13,1) -- cycle;
  \draw[bkTheme, line width=0.8pt, domain=1:13, samples=40, smooth] plot (\x,{1+6.5*exp(-0.3*(\x-1))});
  \or
  \draw[bkGray, line width=0.4pt] (0,4) -- (14,4) (10,-0.5) -- (10,8.5);
  \foreach \kixfx/\kixfy in {3/4,6.5/6.8,6.5/1.2} {\draw[bkRed, line width=1pt] ({\kixfx-1},{\kixfy-1}) -- ({\kixfx+1},{\kixfy+1}) ({\kixfx-1},{\kixfy+1}) -- ({\kixfx+1},{\kixfy-1});}
  \or
  \draw[bkInk, line width=0.6pt, -{Stealth[length=1.4mm]}] (0,4) -- (4,4);
  \draw[bkTheme, line width=0.8pt, fill=bkTheme!12] (4.4,1) rectangle (9.6,7);
  \node[font=\footnotesize, text=bkInk, inner sep=0pt] at (7,4) {$G$};
  \draw[bkInk, line width=0.6pt, -{Stealth[length=1.4mm]}] (10,4) -- (14,4);
  \or
  \draw[bkInk, line width=0.6pt] (1.6,0.5) -- (0.6,0.5) -- (0.6,7.5) -- (1.6,7.5) (5.4,0.5) -- (6.4,0.5) -- (6.4,7.5) -- (5.4,7.5);
  \foreach \kixfx/\kixfy in {2.3/2.3,4.7/2.3,2.3/5.7,4.7/5.7} {\fill[bkInk] (\kixfx,\kixfy) circle (0.5);}
  \draw[bkTheme, line width=0.8pt, domain=0:600, samples=80, smooth] plot ({10.5+3.4*exp(-0.0035*\x)*cos(\x)},{4+3.4*exp(-0.0035*\x)*sin(\x)});
  \fi
\end{tikzpicture}}
\makeatother
%% ===== 表計算・2つの世界の図の型（tools/kata9.md、組 V9-1）。頭文字 kixf =====
\makeatletter
\colorlet{kixfN}{bkOrange!16}% 1行め（成分の番号・周波数）のマス
\colorlet{kixfK}{bkGreen!16}% 2行め（成分の大きさ）のマス
\colorlet{kixfS}{bkYellow!45}% 合計・注目するマス
\newcommand\kixffx[2]{\tikz[baseline=(f.base)]{\node[draw=bkLine, line width=0.6pt, fill=white, minimum height=5.5mm, inner xsep=2mm, font=\sffamily\footnotesize, minimum width=12mm] (c) {#1};
  \node[anchor=west, draw=bkLine, line width=0.6pt, fill=white, minimum height=5.5mm, inner xsep=2mm, minimum width=78mm, text width=74mm, align=left, font=\ttfamily\footnotesize] (f) at ([xshift=1mm]c.east) {#2};}}
\colorlet{kixfW}{bkBlue!80!black}% 時間の世界の線
\colorlet{kixfF}{bkOrange!85!black}% 周波数の世界の棒
\colorlet{kixfP}{bkRed!80!black}% s の世界の点（×）
\pgfplotsset{
  kixf base/.style={bkplot, scale only axis, width=47mm, height=28mm, axis lines=left, grid=none,
    every axis x label/.style={at={(axis description cs:0.5,0)}, anchor=north, yshift=-5mm, font=\sffamily\small, text=bkInk},
    every axis y label/.style={at={(axis description cs:0,1)}, anchor=south, yshift=1.5mm, font=\sffamily\small, text=bkInk}},
  kixf wave/.style={kixf base, extra y ticks={0}, extra y tick labels={0}, extra y tick style={grid=major, major grid style={bkInk!45, line width=0.5pt}}},
  kixf spec/.style={kixf base, ymin=0},
  kixf splane/.style={bkplot, scale only axis, width=47mm, height=47mm, axis equal image, grid=major,
    every axis x label/.style={at={(ticklabel* cs:1)}, anchor=north east, yshift=-0.5mm, font=\sffamily\small, text=bkInk},
    every axis y label/.style={at={(ticklabel* cs:1)}, anchor=south, yshift=1mm, font=\sffamily\small, text=bkInk}},
  kixf sig/.style={kixfW, line width=1.4pt, no markers, samples=241, smooth},
  kixf part/.style={bkGreen!45!black, line width=0.8pt, dash pattern=on 3pt off 1.5pt, no markers, samples=241, smooth},
  kixf bar/.style={ycomb, draw=kixfF, line width=3pt, no markers},
  kixf dot/.style={only marks, mark=*, mark size=1.3pt, kixfW}}
\tikzset{
  kixf world/.style={draw, line width=1pt, rounded corners=1.5mm, fill=white, font=\sffamily\bfseries\small, text=bkInk,
    minimum width=30mm, minimum height=7mm, inner xsep=2mm},
  kixf bridge/.style={-{Stealth[length=2.4mm]}, line width=1pt, bkInk!80},
  kixf blab/.style={font=\sffamily\footnotesize, text=bkInk, inner sep=0.8mm}}
\newcommand\kixfworlds[5][42mm]{%
  \node[kixf world, draw=kixfW] (kixfL) at (23.5mm,#1) {#2};
  \node[kixf world, draw=kixfF] (kixfR) at (99.5mm,#1) {#3};
  \draw[kixf bridge] ([yshift=1.3mm, xshift=1.5mm]kixfL.east) -- node[kixf blab, above] {#4} ([yshift=1.3mm, xshift=-1.5mm]kixfR.west);
  \draw[kixf bridge] ([yshift=-1.3mm, xshift=-1.5mm]kixfR.west) -- node[kixf blab, below] {#5} ([yshift=-1.3mm, xshift=1.5mm]kixfL.east);}
\newcommand\kixfpole[3][kixfP]{\draw[#1, line width=1.3pt] (axis cs:#2,#3) ++(-1.3mm,-1.3mm) -- ++(2.6mm,2.6mm) (axis cs:#2,#3) ++(-1.3mm,1.3mm) -- ++(2.6mm,-2.6mm);}
\newcommand\kixfzero[3][kixfP]{\draw[#1, line width=1.1pt, fill=white] (axis cs:#2,#3) circle[radius=1.2mm];}
\makeatother
%% ===== この章の部品 =====
\makeatletter
% 小見出し（太字の1行）
\newcommand\kixfhead[1]{\par\addvspace{6pt}\Needspace{6\baselineskip}\noindent{\sffamily\bfseries\color{\bkTT}#1}\par\nobreak\vspace{2pt}\noindent\ignorespaces}
% 囲みの中の小見出し（それだけで1行）
\newcommand\kixfbox[1]{\par\addvspace{3pt}\noindent\textbf{#1}\par\vspace{3pt}\noindent\ignorespaces}
% 式を、まん中に1行で置く
\newcommand\kixfline[1]{\par\addvspace{3pt}{\centering #1\par}\addvspace{3pt}\noindent\ignorespaces}
% 大事な式のわく  \kixfframe[はば]{中身}
\newcommand\kixfframe[2][100mm]{\par\addvspace{4pt}{\centering\tikz\node[draw=\bkTT, line width=0.9pt, fill=bkTheme!6, rounded corners=1.5mm, inner xsep=6mm, inner ysep=2.5mm, align=center]{#2};\par}\addvspace{4pt}\noindent\ignorespaces}
% 例題の中の式の変形の表（bkhenkei と同じ形で、行の間を少しつめたもの）
\newenvironment{kixfhenkei}[1][0.48]{\par\addvspace{2pt}\noindent\renewcommand\arraystretch{1.3}%
  \begin{tabular}{@{\hspace{1\zw}}>{$\displaystyle}m{#1\linewidth}<{$}@{\hspace{2\zw}}>{\small\color{bkGray}\raggedright\arraybackslash}m{\dimexpr0.90\linewidth-#1\linewidth\relax}@{}}}%
  {\end{tabular}\par\addvspace{2pt}}
% 例題の途中式の2行め以降
\newcommand\kixfa{\par\vspace{2pt}\noindent\rightskip=0pt plus 8em\relax\hangindent=2.6\zw\hangafter=0\relax}
\newcommand\kixfs{\rightskip=0pt plus 8em\relax}
\makeatother
%% 9-5 微分をかけ算に変える翻訳機：ラプラス変換（第15回：組 V9-4。ページ制限なし。偶数ページ）
%% 「名前と操作」の型（tools/kata9-int.md）で書く。場合分けの表9-6 は A・B・E の行だけ（全部は 9-6）。
%% ほかの巻との約束：5-14:644「e^{-st} をかけて積分する計算」、5-16:605「微分が s のかけ算に変わる。そのしくみが、部分積分」、
%%   5-18:237「e^{-t} がかかると t^n がかかっていても0に近づく」、6-12:238「同じ s」、6-12:899（ヘビサイド）の続き、
%%   6-15「微分を s のかけ算に変える翻訳機」、3-8（部分分数）、冷めるお湯 20+60e^{-0.05t}（10分 56.39 ℃、オイラー法 55.92 ℃）、
%%   RC 10 V・20 kΩ・100 µF・τ=2 s（2秒 6.321 V）、車A s^2+3s+2。数はすべて Python で検算した（scratchpad/v9-4/check1.py・check2.py）。

\bkchapter{9}{5}{微分をかけ算に変える翻訳機：\\ラプラス変換}{\bklevel{5}}{微分をかけ算に変える翻訳機：ラプラス変換}
\bkmichishirube{3-8,3-11,3-12,5-14,5-16,5-18,6-4,6-10,6-12,6-15}{9-6,9-7,9-9}

\begin{goal}
\begin{itemize}[label={\color{\bkTT}●}, leftmargin=1.6\zw, itemsep=0pt, topsep=1pt]
\item \bklevel{3}\enspace 「微分→$s$ をかける」の辞書を使って、冷めるお湯の微分方程式を $s$ の式に書きかえ、表でもどして答えを出せる。
\item \bklevel{4}\enspace ラプラス変換の記号を「名前」として読み、辞書（変換表）の行を、定義の積分で1回だけ作れる。
\item \bklevel{5}\enspace 微分の行 $sF(s)-f(0)$ を部分積分で導き、場合分けの表9-6（A・B・E）を使って微分方程式を変換できる。
\end{itemize}
\end{goal}

\begin{hatena}
\noindent\begin{minipage}[c]{0.6\linewidth}
\raggedright
80\,℃のお湯を、20\,℃の部屋に置きます。温度は、部屋との差の0.05倍の速さ（1分あたり）で下がります（6-1・6-4 の冷めるお湯）。\par\vspace{4pt}
\textbf{10分後の温度を、微分も積分もしないで、「辞書」を引くだけで求められるでしょうか。}
\end{minipage}\hfill
\begin{minipage}[c]{0.36\linewidth}
\centering
\begin{tikzpicture}[x=1mm, y=1mm]
  % カップ
  \filldraw[fill=bkOrange!18, draw=bkInk, line width=0.8pt] (0,0) -- (2,16) -- (18,16) -- (20,0) -- cycle;
  \draw[bkInk, line width=0.8pt] (20,12) arc (90:-90:3.5 and 4.5);
  \foreach \kixfx in {6,10,14} {\draw[bkGray, line width=0.7pt] (\kixfx,18) .. controls (\kixfx-1.5,20.5) and (\kixfx+1.5,22) .. (\kixfx,24.5);}
  \node[font=\footnotesize, anchor=north] at (10,-1.5) {80\,℃};
  \draw[bkInk, line width=0.7pt, -{Stealth[length=2mm]}] (27,8) -- (36,8);
  \node[font=\footnotesize, anchor=south] at (31.5,9) {10分};
  \node[font=\sffamily\bfseries\Large, text=\bkTT] at (42,8) {?};
  % 辞書
  \filldraw[fill=bkTheme!12, draw=\bkTT, line width=0.8pt, rounded corners=0.8mm] (12,-17) rectangle (34,-6);
  \draw[\bkTT, line width=0.8pt] (23,-17) -- (23,-6);
  \node[font=\footnotesize, text=bkInk] at (17.5,-11.5) {$t$};
  \node[font=\footnotesize, text=bkInk] at (28.5,-11.5) {$s$};
  \node[font=\footnotesize, anchor=north] at (23,-18) {辞書（変換表）};
\end{tikzpicture}
\end{minipage}
\end{hatena}

\noindent 6-4 では、冷めるお湯の式を、変数を分けて積分して解きました。6-12 では、ばねとダンパーの式に $e^{st}$ を入れて、特性方程式という $s$ の2次方程式にしました。どちらも、そのたびに頭をひねる必要がありました。

この章では、微分方程式を「\textbf{$s$ の世界}」の式に\textbf{翻訳}する道具、\textbf{ラプラス変換}を学びます。翻訳すると、微分が「$s$ をかける」に変わり、微分方程式が、かけ算・わり算・たし算だけのふつうの方程式になります。答えは、辞書を逆に引いて時間の世界にもどします。制御工学・電気回路で毎日のように使われる道具です。

翻訳の辞書（\textbf{変換表}）は、一度だけ自分で作ります。\bkidx{へんかんひょう}{変換表}作り方は、5-14 の $e$ の積分と、5-16 の部分積分です。新しい考えが多い章です。1回で全部わからなくても大丈夫。辞書を使うやってみようから始めて、あとで「なぜその辞書でよいのか」を確かめます。

\begin{omoidasou}
\bkrecall{6-4}{冷めるお湯 $\tfrac{\dd T}{\dd t}=-0.05(T-20)$、$T(0)=80$ の答えは $T=20+60e^{-0.05t}$（$t$ は分）。10分後 56.39\,℃。}
\bkrecall{3-8}{部分分数：$\tfrac{1}{(s+1)(s+2)}=\tfrac{1}{s+1}-\tfrac{1}{s+2}$。方法2：分母をはらってから、分母を0にする数を代入する。}
\bkrecall{5-14}{$(e^{at})'=ae^{at}$、$\displaystyle\int e^{at}\,\dd t=\frac{1}{a}e^{at}$（$a\neq0$）。$e^{-t}$ は $t$ が大きくなると0に近づく。}
\bkrecall{5-16}{部分積分 $\displaystyle\int f\,g'\,\dd t=f\,g-\int f'\,g\,\dd t$。定積分では $fg$ にもはんいをつける。}
\bkrecall{5-18}{$t^3e^{-t}$ は $t=10$ で0.045、$t=20$ で0.00002。$e^{-t}$ がかかると、$t^n$ がかかっていても0に近づく。}
\bkrecall{6-12}{$m\ddot x+c\dot x+kx=0$ に $e^{st}$ を入れると、特性方程式 $ms^2+cs+k=0$。車A は $s^2+3s+2=0$ で $s=-1,\ -2$。}
\end{omoidasou}

%% ---------------- やってみよう・思い出そう ----------------
\begin{yattemiyou}[辞書を引いて、お湯の式を書きかえよう]
下の辞書を使います。左の列（時間の世界の形）を見つけたら、右の列（$s$ の世界の形）に書きかえます。$T(t)$ を書きかえたものは、大文字のまま $T(s)$ と書きます。
\begin{bkcenter}
\begin{bktable}{colspec={Q[c,m,36mm] Q[c,m,36mm] Q[l,m,42mm]}, row{2-Z}={font=\small}, rowsep=2pt}
時間の世界 & $s$ の世界 & 意味 \\
$T(t)$ & $T(s)$ & 名前を大文字に \\
$\dfrac{\dd T}{\dd t}$ & $sT(s)-T(0)$ & 微分は「$s$ をかけて、はじめの値をひく」 \\
定数 $c$ & $\dfrac{c}{s}$ & 定数は $s$ でわる \\
$e^{at}$ & $\dfrac{1}{s-a}$ & $e$ の肩の数が、分母に入る \\
$f+g$、$cf$ & $F+G$、$cF$ & たし算と定数倍は、そのまま \\
\end{bktable}
\bktabcaption{翻訳の辞書（変換表）のはじめの5行。どの行も、この章の数学の道具で作る。}
\label{bk:9-5:tab-jisho}
\end{bkcenter}
\begin{enumerate}[label={\bkcircled[bkGreen]{\arabic*}}, leftmargin=2.4\zw, itemsep=1pt, topsep=2pt]
\item お湯の式の右辺 $-0.05(T-20)$ を、かっこを外して $-0.05T+1$ にする（$0.05\times20=1$）。
\item 左辺の微分を、辞書で $sT(s)-80$ に（はじめの温度 $T(0)=80$）。右辺の $-0.05T$ は $-0.05T(s)$、定数1は $\tfrac{1}{s}$ に。
\item できた式を、$T(s)$ について、ふつうの1次方程式として解く。
\item 3-8 の部分分数で、辞書にある形（$c/s$ と $c/(s-a)$）に分ける。
\item 辞書を右から左へ引いて、時間の世界にもどす。$t=10$ を入れて、10分後の温度を求める。
\end{enumerate}
\end{yattemiyou}

\noindent ①〜⑤を1行ずつ進めます。②の式は、次のようになります。
\begin{bkhenkei}[0.5]
\frac{\dd T}{\dd t}=-0.05\,T+1 & ①：$-0.05\times(-20)=1$ \\
sT(s)-80=-0.05\,T(s)+\frac{1}{s} & ②：辞書で1項ずつ書きかえる \\
sT(s)+0.05\,T(s)=80+\frac{1}{s} & ③：$T(s)$ の項を左へ \\
(s+0.05)\,T(s)=80+\frac{1}{s} & $T(s)$ でくくる \\
T(s)=\frac{80}{s+0.05}+\frac{1}{s(s+0.05)} & 両辺を $s+0.05$ でわる \\
\end{bkhenkei}
微分が消えて、$s$ のわり算だけの式になりました。2つめの分数は、辞書にない形です。④で、3-8 の方法2（分母をはらって、数を代入する）を使って分けます。
\begin{bkhenkei}[0.5]
\frac{1}{s(s+0.05)}=\frac{A}{s}+\frac{B}{s+0.05} & 分けた形を置く \\
1=A(s+0.05)+Bs & 両辺に $s(s+0.05)$ をかける \\
1=0.05A,\quad A=20 & $s=0$ を入れる（$B$ の項が消える） \\
1=-0.05B,\quad B=-20 & $s=-0.05$ を入れる（$A$ の項が消える） \\
T(s)=\frac{20}{s}+\frac{80-20}{s+0.05} & $A$、$B$ を入れ、分母が同じ項をまとめる \\
=\frac{20}{s}+\frac{60}{s+0.05} & $80-20=60$ \\
\end{bkhenkei}
⑤：辞書を右から左へ引きます。1つめの項は「定数20」です。2つめの項は、$e^{at}$ の行で $a=-0.05$ としたものの60倍なので、「$60e^{-0.05t}$」です。だから
\kixfframe[96mm]{\Large $T=20+60\,e^{-0.05t}$}
6-4 で求めた答えと、まったく同じです（9-0 で予告した翻訳の道すじも、この式でした）。10分後は、次のとおりです。
\kixfline{$T(10)=20+60e^{-0.5}=20+60\times0.6065=56.39$\,℃}
微分も積分もせずに、辞書と、たし算・かけ算・わり算だけで答えが出ました。\bkhatenaref の答えは「求められる。10分後は 56.39\,℃」です。

\begin{bkfigure}
\begin{tikzpicture}
\begin{axis}[bkplot, width=112mm, height=50mm, xmin=0, xmax=32, ymin=0, ymax=90, xtick={0,5,...,30}, ytick={20,40,60,80}, xlabel={$t$（分）}, ylabel={温度（℃）}]
  \addplot[bkthin, dashed, domain=0:31] {20};
  \addplot[bkcurve, domain=0:31, samples=80] {20+60*exp(-0.05*x)};
  \addplot[only marks, mark=square*, mark size=1.6pt, bkOrange] coordinates {(0,80) (5,66.43) (10,55.92) (15,47.80) (20,41.51) (25,36.64) (30,32.88)};
  \node[font=\sffamily\footnotesize, anchor=west, fill=white, inner sep=1pt] at (axis cs:13,78) {青の線：辞書で出した答え $20+60e^{-0.05t}$};
  \node[font=\sffamily\footnotesize, anchor=west, fill=white, inner sep=1pt] at (axis cs:13,68) {オレンジの四角：オイラー法（1分きざみ）};
  \node[font=\sffamily\footnotesize, anchor=north east] at (axis cs:30,17) {部屋の温度 20\,℃};
\end{axis}
\end{tikzpicture}
\bkfigcaption{辞書で出した答え（青の線）と、6-2 のオイラー法（1分きざみ、オレンジの四角）。10分後は 56.39\,℃ と 55.92\,℃。オイラー法は、きざみを細かくすると線に近づく。}
\label{bk:9-5:fig-oyu}
\end{bkfigure}


%% ---------------- 数学の道具 ----------------
\Needspace{8\baselineskip}\bksection{doguu}{数学の道具}

\bkstage{me}{目でわかる}{翻訳して、解いて、もどす}{\bklevel{3}}

\bksubsection{1}{2つの世界と、3つの矢印}

\label{bk:9-5:sekai}やってみようでしたのは、むずかしい計算を、計算しやすい世界に\textbf{翻訳}して解き、もとにもどす、ということです。この考えは、前にも使いました。
3-12 の対数を思い出しましょう。かけ算 $2\times3$ を、対数の表で「たし算」に翻訳し、たして、表を逆に引いてもどしました。計算尺は、このしくみの道具でした。ラプラス変換も、同じ形の考えです。
\begin{bkcenter}
\begin{bktable}{colspec={Q[l,m,30mm] Q[l,m,44mm] Q[l,m,44mm]}, row{2-Z}={font=\small}, rowsep=2pt}
 & 対数（3-12） & ラプラス変換（この章） \\
むずかしい計算 & かけ算 & 微分（微分方程式） \\
翻訳すると & たし算になる & $s$ のかけ算になる \\
翻訳の辞書 & 対数の表 & 変換表 \\
もどし方 & 表を逆に引く & 表を逆に引く（9-7） \\
\end{bktable}
\bktabcaption{対数とラプラス変換は、同じ形の「翻訳の道具」。}
\label{bk:9-5:tab-taisu}
\end{bkcenter}

やってみようでしたことを、1枚の図にまとめます。左の列が\textbf{時間の世界}（$t$ で書いた式とグラフ）、右の列が\textbf{$s$ の世界}（$s$ で書いた式）です。

\begin{bkfigure}
\begin{tikzpicture}[x=1mm, y=1mm, kixfw/.style={draw=\bkTT, line width=0.9pt, rounded corners=1.5mm, fill=white, align=flush center, text width=50mm, inner sep=2mm, font=\small, text=bkInk, minimum height=17mm}]
  \node[font=\sffamily\bfseries, text=\bkTT] at (27,50) {時間の世界（$t$）};
  \node[font=\sffamily\bfseries, text=\bkTT] at (105,50) {$s$ の世界};
  \node[kixfw] (a) at (27,36) {微分方程式\\[1pt]$\dfrac{\dd T}{\dd t}=-0.05\,T+1$、$T(0)=80$};
  \node[kixfw, fill=bkTheme!6] (b) at (105,36) {かけ算・たし算の式\\[1pt]$sT(s)-80=-0.05\,T(s)+\dfrac{1}{s}$};
  \node[kixfw, fill=bkTheme!6] (c) at (105,2) {$s$ の式の答え\\[1pt]$T(s)=\dfrac{20}{s}+\dfrac{60}{s+0.05}$};
  \node[kixfw] (d) at (27,2) {時間の答え\\[1pt]$T=20+60\,e^{-0.05t}$};
  \draw[bk link] (a) -- node[font=\footnotesize, above=0.8mm, align=center] {①翻訳\\$\mathcal L$} (b);
  \draw[bk link] (b) -- node[font=\footnotesize, right=1mm, align=left] {②ふつうの\\方程式で解く} (c);
  \draw[bk link] (c) -- node[font=\footnotesize, below=0.8mm, align=center] {③表で\\もどす} (d);
  \draw[bkInk!60, line width=0.9pt, dash pattern=on 3pt off 2pt, -{Stealth[length=2.2mm]}] (a) -- node[font=\footnotesize, left=1mm, align=right, text=bkInk] {6-4 の\\解き方} (d);
\end{tikzpicture}
\bkfigcaption{翻訳の道すじ。時間の世界で直接解く（点線、6-4）かわりに、$s$ の世界に翻訳し（①）、ふつうの方程式として解き（②）、辞書を逆に引いてもどす（③）。}
\label{bk:9-5:fig-sekai}
\end{bkfigure}

\noindent 遠回りに見えますが、①と③は辞書を引くだけ、②はかけ算・わり算・たし算だけです。むずかしい所が1つもない道に取りかえたのです。

\Needspace{12\baselineskip}
\begin{seikaku}
\kixfbox{表を逆に引いてよいわけ：もどし方は1つに決まる}
③で「辞書を逆に引いてもどす」ことができるのは、\textbf{ちがう関数は、ちがう $F(s)$ に翻訳される}からです。もし、2つのちがう関数 $f(t)$ と $g(t)$ が同じ $F(s)$ になってしまうと、$F(s)$ からもどすときに、どちらにもどせばよいか決まりません。数学では、$t\geqq0$ でとちゅうでとばずにつながった（連続な）関数どうしなら、ある数より大きいすべての $s$ で $F(s)$ が同じとき、$f(t)$ と $g(t)$ も同じであることが証明されています（レルヒの定理）。だから、$\tfrac{60}{s+0.05}$ を見たら、もとは $60e^{-0.05t}$ だと決めてよいのです。ただし、$t<0$ の値は翻訳に入っていないので、もどした答えは $t\geqq0$ の所だけの答えです。
\end{seikaku}

\bksubsection{2}{はじめの値が、自動で入る}

\label{bk:9-5:jidou}6-12 で $e^{st}$ を入れたときは、まず特性方程式で $s$ を決め、そのあとで、はじめの位置や速度に合うように $A$、$B$ を別に決めました。辞書の微分の行 $sT(s)-T(0)$ には、はじめの値 $T(0)$ が入っています。だから、翻訳した式には、はじめから $80$ が入っていて、$T(s)$ を求めると、はじめの値に合った答えが1回で出ます。これが、ラプラス変換のうれしい所の1つです。

\bkstage{te}{手でわかる}{辞書の行を、自分で作る}{\bklevel{4}}

\bksubsection{3}{名前の段：ラプラス変換の記号を読む}

\label{bk:9-5:namae}辞書の行は、どこから来たのでしょう。ラプラス変換は、次の式で決めます。
\kixfframe[112mm]{\Large $\displaystyle\mathcal L[f(t)]=F(s)=\int_0^{\infty}f(t)\,e^{-st}\,\dd t$\\[3pt]\normalsize\sffamily 「$f(t)$ に $e^{-st}$ をかけて、$t=0$ から、どこまでもたし合わせる」}
$\mathcal L$ は、ラプラス（Laplace）の頭文字を、筆記体ふうにした字で、「ラプラス」と読みます（「エル」と読む人もいます）。これを $f(t)$ の\textbf{ラプラス変換}（らぷらすへんかん）といいます。\bkidx{らぷらすへんかん}{ラプラス変換}時間の関数は小文字 $f(t)$、変換したものは大文字 $F(s)$ で書く約束です（$x\to X$、$y\to Y$）。

\kixfdanhead{1}{記号を読む：何に・何をかけて・どこからどこまで}
\begin{enumerate}[label={\bkcircled[\bkTT]{\arabic*}}, leftmargin=2.4\zw, itemsep=2pt, topsep=2pt]
\item \textbf{何に}：時間の関数 $f(t)$（お湯の温度、コンデンサーの電圧、車の位置など）。
\item \textbf{何をかけて}：$e^{-st}$。$s$ は、自分で選ぶ数です（今は正の数と考えます。複素数のときは⑧で）。
\item \textbf{どこからどこまで}：$t=0$ から $\infty$ まで。$t<0$ の値は使いません。スイッチを入れた時刻を0にして、そこから先だけを見ます。
\end{enumerate}
$s$ の値を1つ決めると、積分の値が1つ決まります。だから、変換したものは $s$ の関数 $F(s)$ です。$t$ の関数を、$s$ の関数に取りかえる操作、と言うこともできます。

\kixfhead{記号は「操作の名前」}
第8巻の積分の記号と同じく、$\mathcal L[f(t)]$ は操作の\textbf{名前}です。名前を見ただけでは、まだ数は出ません。実際の値は、次の2つの道具のどちらかに直して出します。
\[ \kixfNt{\mathcal L[f(t)]}=\kixfSt{1}{\int_0^{\infty}f(t)\,e^{-st}\,\dd t} \]
\[ \kixfNt{\mathcal L[e^{-0.05t}]}=\kixfSt{2}{\frac{1}{s+0.05}} \]
上の式は、辞書の行を\textbf{1回だけ作る}ときの直し方（ふつうの定積分）、下の式は、できた辞書を\textbf{使う}ときの直し方（表を引く）です。このように、名前を、計算できる形に直してする計算を\textbf{実際の操作}といいます。ラプラス変換では、ほとんどいつも「表を引く」で終わります。

\kixfhead{$\infty$ までの積分}
$\int_0^{\infty}$ は、「$0$ から $b$ までのふつうの定積分を計算して、$b$ をどこまでも大きくした先」という意味です（8-6 のガウス積分の「範囲をどんどん広げた先」と同じ）。このような積分を\textbf{広義積分}（こうぎせきぶん）とよびます。\bkidx{こうぎせきぶん}{広義積分}先が決まった数に近づかないときは、「変換できない（その $s$ では積分が有限にならない）」と言います。

\bksubsection{4}{なぜ $e^{-st}$ をかけるのか}

\label{bk:9-5:naze}理由は2つあります。1つめは、\textbf{どこまでもたしても、有限の値にするため}です。$f(t)=1$ を、$0$ から $\infty$ まで、そのまま積分すると、面積はいくらでも大きくなります。$e^{-st}$ をかけると、大きな $t$ の所が小さくおさえられ、面積が決まった数になります。

\begin{bkfigure}
\begin{tikzpicture}
\begin{axis}[bkplot, width=112mm, height=60mm, xmin=0, xmax=4.2, ymin=0, ymax=1.3, xtick={0,1,2,3,4}, ytick={0.5,1}, xlabel={$t$}, ylabel={$1\times e^{-st}$}]
  \addplot[draw=none, fill=bkBlue!18, domain=0:4, samples=60] {exp(-2*x)} \closedcycle;
  \addplot[bkthin, domain=0:4] {1};
  \addplot[bkcurve2, dashed, domain=0:4, samples=60] {exp(-x)};
  \addplot[bkcurve, domain=0:4, samples=60] {exp(-2*x)};
  \addplot[bkcurve3, dash pattern=on 1pt off 1.5pt, domain=0:4, samples=60] {exp(-4*x)};
  \node[font=\sffamily\footnotesize, anchor=south, fill=white, inner sep=1pt] at (axis cs:2.6,1.03) {灰色の線：$f(t)=1$ のまま（面積は限りなく大きい）};
  \node[font=\sffamily\footnotesize, anchor=west, fill=white, inner sep=1pt] at (axis cs:1.6,0.84) {オレンジの破線：$s=1$、$e^{-t}$、面積1};
  \node[font=\sffamily\footnotesize, anchor=west, fill=white, inner sep=1pt] at (axis cs:1.6,0.69) {青の線：$s=2$、$e^{-2t}$、面積 0.5};
  \node[font=\sffamily\footnotesize, anchor=west, fill=white, inner sep=1pt] at (axis cs:1.6,0.54) {緑の点線：$s=4$、$e^{-4t}$、面積 0.25};
\end{axis}
\end{tikzpicture}
\bkfigcaption{$f(t)=1$ に $e^{-st}$ をかけたグラフ。$s$ が大きいほど早く0に近づき、下の面積（$s=2$ の所をぬった）は小さくなる。面積は $\frac{1}{s}$：これが辞書の「定数1 → $\frac{1}{s}$」。}
\label{bk:9-5:fig-menseki}
\end{bkfigure}

\noindent 5-18 で見たように、$e^{-t}$ は、$t^3$ のような累乗の関数よりずっと速く0に近づきます。だから、$f(t)$ が $t$、$t^2$、$e^{0.5t}$ のようにふえる関数でも、$s$ を十分大きくすれば、$f(t)e^{-st}$ は0に近づき、面積が決まった数になります。

2つめの理由は、\textbf{微分が $s$ のかけ算に変わるため}です。これは⑥の部分積分で見ます。$e^{-st}$ は微分しても $-s$ 倍になるだけの、形の変わらない関数（5-14）なので、部分積分をすると $s$ だけが前に出てくるのです。

\bksubsection{5}{直す段：辞書の行を、定義の積分で作る}

\label{bk:9-5:gyou}\kixfdanhead{2}{ふつうの定積分に直す（1回だけ）}
辞書の行を、定義の積分で1行ずつ作ります。どの行も、$0$ から $b$ までの定積分を計算して、$b$ を大きくします。

\kixfhead{定数1の行：$\mathcal L[1]=\dfrac{1}{s}$}
\begin{bkhenkei}[0.5]
\int_0^{b}1\times e^{-st}\,\dd t=\Bigl[-\frac{1}{s}e^{-st}\Bigr]_0^{b} & $e^{-st}$ の積分（5-14。$a=-s$） \\
=-\frac{1}{s}e^{-sb}+\frac{1}{s} & 上のはし $b$ と下のはし0を入れて、ひく \\
\to0+\frac{1}{s}=\frac{1}{s} & $b$ を大きくすると $e^{-sb}\to0$（$s>0$ のとき） \\
\end{bkhenkei}
定数 $c$ なら、$c$ を前に出して $\tfrac{c}{s}$ です（定数倍は積分の外に出せる：5-11）。$s$ が0や負の数だと、$e^{-sb}$ は0に近づかず、積分は有限になりません。だから、この行は「$s>0$ のとき」の行です。

\kixfhead{表計算で確かめる}
定義の積分を、表計算の台形の和（5-10）で近い値にしてみます。$f(t)=1$、$s=2$ として、$t=0$ から10まで $\Delta t=0.1$ ずつ、$f$ と $e^{-2t}$ とその積を並べ、となりどうしの台形（上の値と下の値の平均×0.1）をたします。$\infty$ のかわりに $t=10$ で止めてよいのは、$e^{-20}=2.1\times10^{-9}$ で、その先はほとんど0だからです。
\begin{bkcenter}
\kixffx{E3}{=(D2+D3)/2*0.1}\par\vspace{1mm}
\bksheet{5}{17}{%
1 & \SetCell{c} $t$ & \SetCell{c} $f(t)$ & \SetCell{c} $e^{-2t}$ & \SetCell{c} $f\,e^{-2t}$ & \SetCell{c} 台形 \\
2 & 0 & 1 & 1.0000 & 1.0000 & \\
3 & 0.1 & 1 & 0.8187 & 0.8187 & 0.0909 \\
4 & 0.2 & 1 & 0.6703 & 0.6703 & 0.0745 \\
5 & 0.3 & 1 & 0.5488 & 0.5488 & 0.0610 \\
$\vdots$ & $\vdots$ & $\vdots$ & $\vdots$ & $\vdots$ & $\vdots$ \\
102 & 10 & 1 & 0.0000 & 0.0000 & 0.0000 \\
103 & \SetCell{c} 合計 & & & & \SetCell{bg=kixfS} 0.50167 \\}
\bktabcaption{$\int_0^\infty f\,e^{-2t}\,\dd t$ を台形でたす表（$\Delta t=0.1$、$t=10$ まで）。C2 は \texttt{=EXP(-2*A2)}、D2 は \texttt{=B2*C2}、E103 は \texttt{=SUM(E3:E102)}。表の値は $\frac1s=0.5$。}
\label{bk:9-5:tab-sheet}
\end{bkcenter}
表計算の値 0.50167 は、辞書の $\tfrac{1}{s}=\tfrac{1}{2}=0.5$ に近い値です。$\Delta t=0.01$ にすると 0.50002 になり、さらに近づきます。辞書の行は、この「たし合わせ」を、式で一度に計算したものなのです。

\kixfhead{ステップ関数 $H(t)$：時刻0にスイッチを入れる}
\noindent\begin{minipage}[c]{0.66\linewidth}
\raggedright\setlength{\parindent}{1\zw}
ラプラス変換は $t\geqq0$ だけを見るので、「定数1」は、正確には「$t<0$ では0、$t\geqq0$ では1」という関数のことです。これを\textbf{ステップ関数}（すてっぷかんすう。ヘビサイドの階段関数）といい、$H(t)$ と書きます。\bkidx{すてっぷかんすう}{ステップ関数}「時刻0にスイッチを入れて、電池の1\,V をつなぐ」ようすを表す関数で、$\mathcal L[H(t)]=\tfrac{1}{s}$ です。$H$ は、6-12 のコラムに出てきたヘビサイドの頭文字です。
\end{minipage}\hfill
\begin{minipage}[c]{0.3\linewidth}
\centering
\begin{tikzpicture}[x=1mm, y=1mm]
  \draw[bkInk, line width=0.7pt, -{Stealth[length=2mm]}] (-16,0) -- (18,0) node[right, font=\footnotesize] {$t$};
  \draw[bkInk, line width=0.7pt, -{Stealth[length=2mm]}] (0,-2) -- (0,16) node[above, font=\footnotesize] {$H(t)$};
  \draw[bkBlue, line width=1.6pt] (-15,0) -- (0,0);
  \draw[bkBlue, line width=1.6pt] (0,10) -- (16,10);
  \draw[bkBlue, line width=0.8pt, dash pattern=on 1.5pt off 1pt] (0,0) -- (0,10);
  \fill[bkBlue] (0,10) circle (0.9);
  \node[font=\footnotesize, anchor=east] at (-1.5,10) {1};
  \node[font=\footnotesize, anchor=north east] at (-1,-0.5) {0};
\end{tikzpicture}
\end{minipage}\par\vspace{2pt}

\kixfhead{$e^{at}$ の行：$\mathcal L[e^{at}]=\dfrac{1}{s-a}$}
$e^{at}\times e^{-st}=e^{-(s-a)t}$（肩のたし算：5-14）なので、上の計算の $s$ を $s-a$ に取りかえるだけです。
\begin{bkhenkei}[0.5]
\int_0^{b}e^{at}e^{-st}\,\dd t=\int_0^{b}e^{-(s-a)t}\,\dd t & 肩をまとめる \\
=\Bigl[-\frac{1}{s-a}e^{-(s-a)t}\Bigr]_0^{b} & 5-14（肩の係数 $-(s-a)$ でわる） \\
\to\frac{1}{s-a} & $s>a$ なら、$b$ を大きくすると $e^{-(s-a)b}\to0$ \\
\end{bkhenkei}
お湯の $e^{-0.05t}$ は $a=-0.05$ なので $\tfrac{1}{s+0.05}$ です（$s>-0.05$ のとき）。「$s>a$ のとき」という条件は、「$e^{-st}$ が、$e^{at}$ のふえ方に勝つくらい速く小さくなるとき」という意味です。

\kixfhead{$t$ の行：$\mathcal L[t]=\dfrac{1}{s^2}$（部分積分）}
$t\times e^{-st}$ は、5-16 の部分積分を使います。微分すると簡単になる $t$ を $f$、積分しやすい $e^{-st}$ を $g'$ にします（$g=-\frac{1}{s}e^{-st}$）。
\begin{bkhenkei}[0.64]
\int_0^{\infty}t\,e^{-st}\,\dd t & 求めたいもの \\
=\Bigl[t\cdot\Bigl(-\frac{1}{s}e^{-st}\Bigr)\Bigr]_0^{\infty}-\int_0^{\infty}1\cdot\Bigl(-\frac{1}{s}e^{-st}\Bigr)\dd t & 部分積分 $fg-\int f'g$ \\
=(0-0)+\frac{1}{s}\int_0^{\infty}e^{-st}\,\dd t & はしの項は0（下の説明） \\
=\frac{1}{s}\times\frac{1}{s}=\frac{1}{s^2} & 残りは $\mathcal L[1]=\frac{1}{s}$ \\
\end{bkhenkei}
はしの項 $\Bigl[-\tfrac{t}{s}e^{-st}\Bigr]_0^{\infty}$ は、$t=0$ では $0$、$t$ が大きくなると、$te^{-st}$ が0に近づくので0です（5-18：$e^{-t}$ がかかると、$t^n$ がかかっていても0に近づく）。

\kixfhead{Python で確かめる}
Python では、$\infty$ までの積分を数値で計算する \texttt{quad}（scipy）が使えます。$s=2$ のときの $\mathcal L[1]$・$\mathcal L[t]$・$\mathcal L[e^{-t}]$ を計算し、辞書の値 $\tfrac{1}{s}$・$\tfrac{1}{s^2}$・$\tfrac{1}{s+1}$ とくらべます。
\Needspace{17\baselineskip}
\begin{bkcode}[Python]
import numpy as np
from scipy.integrate import quad
s = 2
fs = [lambda t: 1, lambda t: t, lambda t: np.exp(-t)]
table = [1/s, 1/s**2, 1/(s+1)]      \# 辞書の値
for f, F in zip(fs, table):
    v = quad(lambda t: f(t)*np.exp(-s*t), 0, np.inf)[0]
    print(round(v, 5), round(F, 5))
\end{bkcode}
\begin{bkcodeout}[実行結果]
0.5 0.5
0.25 0.25
0.33333 0.33333
\end{bkcodeout}
左が数値の積分、右が辞書の値です。どれも同じになりました。

\bksubsection{6}{微分の行：$\mathcal L[f']=sF(s)-f(0)$}

\label{bk:9-5:bibun}いよいよ、この章でいちばん大事な行です。5-16 で「$e^{-st}$ をかけて積分すると、微分が $s$ のかけ算に変わります。そのしくみが、部分積分です」と予告しました。$f'(t)$（$f$ を微分したもの）のラプラス変換を、部分積分で計算します。こんどは、$e^{-st}$ を $f$ の役（微分するほう）、$f'(t)$ を $g'$ の役（積分するほう）にします。$f'$ を積分すると $f$ にもどる（5-11）ので、$g=f(t)$ です。
\begin{bkhenkei}[0.52]
\int_0^{\infty}e^{-st}f'(t)\,\dd t & 求めたいもの \\
=\Bigl[e^{-st}f(t)\Bigr]_0^{\infty}-\int_0^{\infty}(-s\,e^{-st})\,f(t)\,\dd t & 部分積分。$(e^{-st})'=-se^{-st}$ \\
=\bigl(0-e^{0}f(0)\bigr)+s\int_0^{\infty}f(t)\,e^{-st}\,\dd t & はしの項（下の説明）。$-s$ を外へ \\
=-f(0)+s\,F(s) & 残った積分は、$f$ のラプラス変換 $F(s)$ \\
\end{bkhenkei}
\kixfframe[100mm]{\Large $\mathcal L[f'(t)]=s\,F(s)-f(0)$}
\begin{itemize}[leftmargin=1.2\zw, itemsep=2pt, topsep=2pt, label={\color{\bkTT}●}]
\item \textbf{$s$ が前に出る}のは、$e^{-st}$ を微分すると $-s$ 倍になるから。部分積分で、微分が $f'$ から $e^{-st}$ に移り、そのとき $s$ が出てきます。
\item \textbf{$-f(0)$ は、はしの項}です。$t=\infty$ の側は、$e^{-st}$ が $f(t)$ のふえ方に勝って0。$t=0$ の側は $e^0f(0)=f(0)$ で、それをひきます。
\end{itemize}
はじめの値 $f(0)$ が、計算のしくみから自然に入ってきます。②で見た「はじめの値が自動で入る」は、この部分積分のはしの項のおかげでした。

\kixfhead{数で確かめる：$f=e^{-2t}$}
$f=e^{-2t}$ なら、$f'=-2e^{-2t}$ です。左辺を辞書で直接、右辺を微分の行で計算して、くらべます。
\begin{bkhenkei}[0.5]
\mathcal L[f']=\mathcal L[-2e^{-2t}]=\frac{-2}{s+2} & 直接：$e^{at}$ の行で $a=-2$、$-2$ 倍 \\
sF(s)-f(0)=\frac{s}{s+2}-1 & 微分の行：$F=\frac{1}{s+2}$、$f(0)=e^0=1$ \\
=\frac{s-(s+2)}{s+2}=\frac{-2}{s+2} & 通分する \\
\end{bkhenkei}
同じになりました。

\bkkokomade[辞書を引いて微分方程式を $s$ の式に書きかえ、部分分数と辞書でもどせること、辞書の行が「$e^{-st}$ をかけて0から $\infty$ までの積分」で作れること、微分の行 $sF-f(0)$ が部分積分から出ることがわかれば十分です。この先は、2回微分の行、$s$ が複素数のとき、場合分けの表です。]

\Needspace{9\baselineskip}\bkstage{shiki}{式でわかる}{2回微分、複素数の $s$、場合分け}{\bklevel{5}}

\bksubsection{7}{2回微分の行と、6-12 の特性方程式}

\label{bk:9-5:nikai}2回微分 $f''$ は、「$f'$ をもう一度微分したもの」です。微分の行を2回使います。
\begin{bkhenkei}[0.52]
\mathcal L[f'']=s\,\mathcal L[f']-f'(0) & 微分の行を、$f$ のかわりに $f'$ に使う \\
=s\bigl(sF(s)-f(0)\bigr)-f'(0) & $\mathcal L[f']=sF-f(0)$ を入れる \\
=s^2F(s)-s\,f(0)-f'(0) & かっこを外す \\
\end{bkhenkei}
\kixfframe[100mm]{\Large $\mathcal L[f''(t)]=s^2F(s)-s\,f(0)-f'(0)$}
はじめの位置 $f(0)$ と、はじめの速度 $f'(0)$ の両方が、自動で入ります。

\kixfhead{車A を翻訳する}
3-10・6-12 の車A $\ddot x+3\dot x+2x=0$（はじめ $x(0)=1$、$\dot x(0)=0$）を翻訳します。
\begin{bkhenkei}[0.52]
\bigl(s^2X-s\cdot1-0\bigr)+3\bigl(sX-1\bigr)+2X=0 & 1項ずつ辞書で（$\ddot x$、$\dot x$、$x$） \\
(s^2+3s+2)\,X=s+3 & $X$ の項を左に、残りを右に \\
X(s)=\frac{s+3}{s^2+3s+2}=\frac{s+3}{(s+1)(s+2)} & わる。分母を因数分解 \\
\end{bkhenkei}
分母の $s^2+3s+2$ は、6-12 で $e^{st}$ を入れて作った特性方程式の左辺と同じです。6-12 で「第9巻のラプラス変換でも同じ $s$ を使います」と書いたのは、このためでした。この $X(s)$ を部分分数に分けて表でもどす方法は 9-7 でくわしく学びますが、やってみようと同じ方法で $\tfrac{2}{s+1}-\tfrac{1}{s+2}$ に分かれ、6-12 と同じ答え $x=2e^{-t}-e^{-2t}$ になります。

\begin{motto}[5]
\kixfbox{6-12 の解き方と、ラプラス変換の解き方をくらべる}
車A（$\ddot x+3\dot x+2x=0$、$x(0)=1$、$\dot x(0)=0$）を、2つの方法で解いた手順を並べます。
\begin{bkcenter}
\begin{bktable}{colspec={Q[c,m,8mm] X[l,m] X[l,m]}, row{2-Z}={font=\small}, rowsep=2pt, colsep=3pt}
 & 6-12（$e^{st}$ を入れる） & 9-5（ラプラス変換） \\
① & $e^{st}$ を入れて、特性方程式 $s^2+3s+2=0$ & 両辺を変換：$(s^2+3s+2)X=s+3$ \\
② & 解 $s=-1$、$-2$ から $x=Ae^{-t}+Be^{-2t}$ & $X=\frac{s+3}{(s+1)(s+2)}$ を部分分数に \\
③ & はじめの値から連立方程式 $A+B=1$、$-A-2B=0$ & $s=-1$、$-2$ を代入して係数 2、$-1$ \\
④ & $A=2$、$B=-1$ & 表でもどす \\
答え & $x=2e^{-t}-e^{-2t}$ & $x=2e^{-t}-e^{-2t}$ \\
\end{bktable}
\end{bkcenter}
6-12 では、はじめの値を最後に連立方程式で入れました。ラプラス変換では、①の翻訳のときに、はじめの値が $s+3$ の中に入っています。式がもっと複雑になったり、右辺に入力があったりしても、ラプラス変換の手順は同じです（9-7）。
\end{motto}

\bksubsection{8}{$s$ は複素数でもよい：$s=\sigma+j\omega$}

\label{bk:9-5:fukuso}ここまで、$s$ は正の数として考えてきました。じつは、$s$ は複素数でもかまいません。6-10⑪で学んだように、$s=\sigma+j\omega$（$\sigma$ はシグマ、$\omega$ はオメガ）なら、
\kixfline{$e^{-st}=e^{-\sigma t}\times e^{-j\omega t}$}
です。$e^{-\sigma t}$ は「小さくする」役（$\sigma>0$ なら、$t$ が大きいほど小さく）、$e^{-j\omega t}$ は「回る」役（長さ1のまま、角度 $-\omega t$ で回る矢印）です。だから、$e^{-st}$ をかけることは、「$\sigma$ でだんだん小さくしながら、$\omega$ の速さで回してかける」ことです。

\begin{bkfigure}
\begin{tikzpicture}
\begin{axis}[bkplot, width=112mm, height=52mm, xmin=0, xmax=6.4, ymin=-1.15, ymax=1.15, xtick=\empty, ytick={-1,1}, xlabel={$t$}]
  \addplot[bkthin, dashed, domain=0:6.2, samples=60] {exp(-0.5*x)};
  \addplot[bkthin, dashed, domain=0:6.2, samples=60] {-exp(-0.5*x)};
  \addplot[bkcurve, domain=0:6.2, samples=200] {exp(-0.5*x)*cos(deg(3*x))};
  \node[font=\sffamily\footnotesize, anchor=west] at (axis cs:2.5,0.85) {灰色の破線：$\pm e^{-\sigma t}$（小さくする役）};
  \node[font=\sffamily\footnotesize, anchor=west] at (axis cs:2.5,-0.85) {青の線：$e^{-\sigma t}\cos\omega t$（回る矢印の影）};
\end{axis}
\end{tikzpicture}
\bkfigcaption{$s=0.5+3j$ のときの $e^{-st}$ の実部 $e^{-0.5t}\cos3t$。回る矢印の影がゆれながら、$e^{-0.5t}$ の中に小さくなっていく。}
\label{bk:9-5:fig-fukuso}
\end{bkfigure}

\kixfhead{周波数の世界は、$s$ の世界の $s=j\omega$ の所}
9-0 で、「周波数の世界は、$s$ の世界の一部」と予告しました。9-4 のフーリエ変換は、$f(t)$ に $e^{-j\omega t}$ をかけて、$-\infty$ から $\infty$ までたしました。$f(t)$ が $t<0$ で0なら、これはラプラス変換で $\sigma=0$（$s=j\omega$）にしたものです（その積分が有限になるとき）。ラプラス変換は、フーリエ変換に「小さくする役」$e^{-\sigma t}$ をつけ足して、ふえていく関数（$t$、$e^{0.5t}$ など）でも積分が有限になるようにしたもの、と見ることもできます。

\kixfhead{$s$ の条件は、$s$ の平面の「右側」}
$\mathcal L[e^{at}]=\tfrac{1}{s-a}$ の計算は、$s$ が複素数でも同じで、条件「$s>a$」は「$s$ の実部 $\sigma$ が $a$ より大きい」に変わります。$s$ の平面（6-9 の複素平面。横が $\sigma$、たてが $\omega$）では、たての線 $\sigma=a$ より右側の全部です（図\ref{bk:9-5:fig-ten}の右の平面で、$\mathcal L[e^{-t}]=\frac{1}{s+1}$ なら、×の $-1$ を通るたての線より右）。



\kixfhead{なぜ、6-12 と同じ $s$ なのか}
$e^{at}$ を変換すると $\tfrac{1}{s-a}$ で、$s=a$ のとき分母が0になります。つまり、\textbf{時間の世界の動き $e^{at}$ の $a$ が、$s$ の世界では「分母が0になる場所」として現れる}のです。車A の $X(s)=\tfrac{s+3}{(s+1)(s+2)}$ の分母が0になる $s=-1$、$-2$ は、6-12 で特性方程式を解いて出した $s$、つまり動き $e^{-t}$、$e^{-2t}$ の肩の数そのものです。だから、6-12 の $s$ と、ラプラス変換の $s$ は同じ文字で書くのです（9-0 で「偶然ではありません」と予告したわけ）。分母が0になるこの場所を、9-9 では\textbf{極}とよび、その場所で動きが「落ち着くか・ゆれるか・あばれるか」を読みます。

\begin{bkfigure}
\begin{tikzpicture}
\kixfworlds[63mm]{時間の世界}{$s$ の世界}{ラプラス変換}{逆変換}
\begin{axis}[kixf wave, at={(0mm,0mm)}, height=47mm, xmin=0, xmax=5, ymin=-0.2, ymax=1.1, xtick={0,1,...,5}, ytick={0.2,0.4,0.6,0.8,1}, yticklabel style={/pgf/number format/fixed}, xlabel={時間 $t$}, ylabel={$x$}]
  \addplot[kixf sig, domain=0:5] {2*exp(-x)-exp(-2*x)};
  \addplot[kixf part, domain=0.61:5] {2*exp(-x)};
  \addplot[kixf part, dash pattern=on 1pt off 1.2pt, domain=0.81:5] {-exp(-2*x)};
  \node[font=\sffamily\footnotesize, anchor=south west, fill=white, inner sep=1pt] at (axis cs:1.5,0.62) {$x=2e^{-t}-e^{-2t}$（太い線）};
  \draw[bkInk!70, line width=0.5pt] (axis cs:3.0,0.12) -- (axis cs:3.2,0.32);
  \node[font=\sffamily\footnotesize, anchor=south west, fill=white, inner sep=1pt] at (axis cs:3.1,0.32) {$2e^{-t}$（破線）};
  \node[font=\sffamily\footnotesize, anchor=north west, fill=white, inner sep=1pt] at (axis cs:1.5,-0.03) {$-e^{-2t}$（点線）};
\end{axis}
\begin{axis}[kixf splane, at={(76mm,0mm)}, xmin=-3, xmax=2, ymin=-2.5, ymax=2.5, xtick={-2,-1}, ytick={-2,-1,1,2}, xticklabel style={yshift=-1.4mm}, xlabel={実部 $\sigma$}, ylabel={虚部 $\omega$}]
  \kixfpole{-1}{0}\kixfpole{-2}{0}
  \node[font=\sffamily\footnotesize, anchor=south, fill=white, inner sep=1pt] at (axis cs:-1.5,0.35) {車A の×};
  \node[font=\sffamily\footnotesize, anchor=west, fill=white, inner sep=1pt, align=left] at (axis cs:0.2,1.4) {たての軸が\\$s=j\omega$ の所\\（周波数の世界）};
\end{axis}
\end{tikzpicture}
\bkfigcaption{車A の動き（左）と、$X(s)$ の分母が0になる場所（右の×）。×の $-1$、$-2$ が、動きの $e^{-t}$、$e^{-2t}$ の肩の数。右のたての軸（$\sigma=0$、$s=j\omega$）が、9-4 の周波数の世界にあたる所。}
\label{bk:9-5:fig-ten}
\end{bkfigure}

\bksubsection{9}{たし算と定数倍は、そのまま（線形）}

\label{bk:9-5:senkei}\kixfdanhead{3}{計算する前に、楽な道をさがす：項ごとに分ける}
ラプラス変換は積分なので、たし算と定数倍はそのまま通ります（5-11 の定積分の性質）。
\kixfframe[104mm]{\Large $\mathcal L[a\,f(t)+b\,g(t)]=a\,F(s)+b\,G(s)$}
5-18 で、「答えどうしをたしても答え」になる式を\textbf{線形}な式とよび\bkidx{せんけい}{線形}、「第6巻・第9巻の中心になる考え」と書きました。ラプラス変換も線形です。だから、いくつかの項の和は、項ごとに辞書を引いてたせばよいのです。たとえば
\kixfline{$\mathcal L[20+60e^{-0.05t}]=20\times\dfrac{1}{s}+60\times\dfrac{1}{s+0.05}=\dfrac{20}{s}+\dfrac{60}{s+0.05}$}
で、やってみようの $T(s)$ にもどります。

ただし、\textbf{かけ算は、そのまま通りません}。$\mathcal L[f\,g]$ は、$F(s)\times G(s)$ ではありません。$f=g=1$ なら、$\mathcal L[1\times1]=\tfrac{1}{s}$ ですが、$\mathcal L[1]\times\mathcal L[1]=\tfrac{1}{s^2}$ です。$s$ の世界のかけ算 $F(s)G(s)$ が時間の世界の何になるかは、9-6 の「たたみこみ」で学びます。

\bksubsection{10}{場合分けの表9-6（はじめの3行）}

\label{bk:9-5:baai}ラプラス変換を求めるとき、「どの場合に、どう直すか」を表にします。この章では、表9-6 のはじめの3行（A・B・E）を使います。9-6 で、ずらす・おくれる・$t$ をかけるなどの行（C・D・F〜I）を足して、表と流れ図を完成させます。

\kixfcasetitle{9-6}{ラプラス変換：どんなときに、どう求める？（はじめの3行）}
\begin{bkcenter}
\begin{bktable}{colspec={Q[c,m,10mm] Q[l,m,33mm] X[l,m] Q[c,m,13mm]}, row{2-Z}={font=\small}, rowsep=3pt, colsep=4pt}
記号 & こんなとき & 手順 & 例題 \\
\kixfcase{A} & 辞書（表）にある形（$1$、$t$、$e^{at}$ など） & \kixfop{2}\enspace 表を引く。まだない行は、定義の積分（\kixfop{1}）で1回だけ作る & 1・7 \\
\kixfcase{B} & 表の形の和・定数倍 & \kixfdan{3}\enspace 線形：項ごとに A で引いて、たす & 7 \\
\kixfcase{E} & $f'$・$f''$ のある式（微分方程式）を変換する & \kixfop{2}\enspace $f'\to sF-f(0)$、$f''\to s^2F-sf(0)-f'(0)$。ほかの項は A・B で & 2〜6 \\
\end{bktable}
\end{bkcenter}

\noindent どの場合かは、図\ref{bk:9-5:fig-flow}の流れ図で、上から順に問いに答えて決めます。9-6 の流れ図は、この図に問いを足したものです。

\kixfhead{どの場合かな？}
図\ref{bk:9-5:fig-flow}の流れ図を上からたどって、先に自分で考えてから答えを見ましょう。
\begin{enumerate}[label={\bkcircled[\bkTT]{\arabic*}}, leftmargin=2.4\zw, itemsep=3pt, topsep=2pt]
\item $\mathcal L[e^{3t}]$\par\nobreak\noindent 答え：\kixfcase{A}。$e^{at}$ の行で $a=3$。$\tfrac{1}{s-3}$（$s>3$）。
\item $\mathcal L[5-2t]$\par\nobreak\noindent 答え：\kixfcase{B}→\kixfcase{A}。2つの項の和。$\tfrac{5}{s}-\tfrac{2}{s^2}$。
\item RC 回路の式 $2\tfrac{\dd V}{\dd t}+V=10$ の変換\par\nobreak\noindent 答え：\kixfcase{E}。$f'$ のある式。
\end{enumerate}

\begin{bkfigure}
\begin{tikzpicture}[x=1mm, y=0.9mm]
  \node[kixf do, text width=76mm] (s) at (46,0) {記号を読む（名前の段）：何を変換するか\\（$t$ の関数か、微分方程式か）。定数倍は前に出す};
  \node[kixf q] (q1) at (33,-22) {$f'$・$f''$ のある式を変換？};
  \node[kixf do, text width=40mm, anchor=west] (rE) at (71,-22) {$f'\to sF-f(0)$ など\\ほかの項は A・B で\enspace\kixfcase{E}};
  \node[kixf q] (q2) at (33,-42) {いくつかの項の和？};
  \node[kixf do, text width=40mm, anchor=west] (rB) at (71,-42) {項に分けて、項ごとに\\下の問いへ\enspace\kixfcase{B}};
  \node[kixf do, text width=40mm] (rA) at (33,-62) {表を引く（ない行は\\定義の積分で作る）\enspace\kixfcase{A}};
  \node[kixf do, text width=104mm] (end) at (60,-80) {確かめ：表計算の数値（$s=2$ など）と同じか};
  \draw[kixf go] (s.south -| q1) -- (q1);
  \foreach \kixfq/\kixfr in {q1/rE,q2/rB} {\draw[kixf go] (\kixfq) -- node[kixf yn, above] {はい} (\kixfr);}
  \foreach \kixfq/\kixfr in {q1/q2,q2/rA} {\draw[kixf go] (\kixfq) -- node[kixf yn, left=1mm] {いいえ} (\kixfr);}
  \draw[kixf go] (rB.south) -- ++(0,-3) -| ([xshift=10mm]rA.north);
  \foreach \kixfn in {rE,rA} {\draw[bkInk!60, line width=0.7pt, shorten <=1.5mm] (\kixfn.east) -- (\kixfn.east -| 127,0);}
  \draw[bkInk!60, line width=0.7pt] (127,-22) -- (127,-80);
  \draw[kixf go, shorten <=0pt] (127,-80) -- (end.east);
\end{tikzpicture}
\bkfigcaption{9-5 で使う流れ図（9-6 の流れ図のはじめの部分）。ひし形は問い、四角はすること。B のあとは、項ごとに左の列の下（A）へ。計算した四角からは、右の線をたどって最後の「確かめ」へ。}
\label{bk:9-5:fig-flow}
\end{bkfigure}



\kixfhead{例題を、流れ図でたどってみよう}
この章の例題1〜7を、図\ref{bk:9-5:fig-flow}の流れ図の上からたどると、次のようになります。じぶんでとこう・章末問題も、同じ番号の例題と同じ道すじです。
\begin{bkcenter}
\begin{bktable}{colspec={Q[c,m,10mm] Q[l,m,38mm] X[l,m] Q[c,m,16mm]}, row{2-Z}={font=\small}, rowsep=0.8pt, colsep=3pt}
例題 & 題 & 「はい」と答えた問い & 着いた所 \\
1 & 辞書の行を作る & なし（どちらも「いいえ」） & \kixfcase{A} \\
2 & 冷めるコーヒー & $f'$・$f''$ のある式を変換？ & \kixfcase{E} \\
3 & RC の充電 & $f'$・$f''$ のある式を変換？ & \kixfcase{E} \\
4 & RL 回路 & $f'$・$f''$ のある式を変換？ & \kixfcase{E} \\
5 & 車A & $f'$・$f''$ のある式を変換？ & \kixfcase{E} \\
6 & 微分の行を確かめる & $f'$・$f''$ のある式を変換？ & \kixfcase{E} \\
7 & たし算の形 & いくつかの項の和？ → 項ごとに A & \kixfcase{B}\kixfcase{A} \\
\end{bktable}
\bktabcaption{例題1〜7の流れ図の道すじ。}
\label{bk:9-5:tab-michi}
\end{bkcenter}

\begin{seikaku}
\kixfbox{どんな関数でも変換できるわけではない}
ラプラス変換ができるのは、$f(t)$ が、ある数 $a$ について $e^{at}$ より速くはふえない（$|f(t)|\leqq Me^{at}$ となる数 $M$ がある）ときです。そのとき、実部が $a$ より大きい $s$ で積分が有限になります。$e^{t^2}$ のように、どんな $e^{at}$ よりも速くふえる関数は、変換できません。工学で出てくる関数（定数・$t^n$・$e^{at}$・$\sin$・$\cos$ とその組み合わせ）は、どれも変換できます。また、$f(t)$ がとちゅうでとんでもかまいません（9-6 のパルス）。
\end{seikaku}



\bkkigocards{\bkkigocard{\scalebox{0.72}{$\mathcal L[f(t)]$}}{ラプラス・エフ・ティー（エルと読む人もいる）}{$f(t)$ に $e^{-st}$ をかけて0から $\infty$ までたしたもの。大文字で $F(s)$ と書く}{Laplace transform}{9-5}}
{\bkkigocard{$H(t)$}{エイチ・ティー（ステップ関数）}{$t<0$ で0、$t\geqq0$ で1。$u(t)$・$1(t)$ とも書く}{step function}{9-5}}

\begin{kotoba}
\small
\bkword{ラプラス変換}{$f(t)$ に $e^{-st}$ をかけて、$t=0$ から $\infty$ までたし合わせたもの $F(s)$。操作の名前。実際は、表を引いて求める。}
\bkword{$s$ の世界}{$s$ で書いた式の世界。微分が $s$ のかけ算になる。$s=\sigma+j\omega$ は複素数でもよい。}
\bkword{変換表（辞書）}{時間の関数と、そのラプラス変換を並べた表。行は、定義の積分で1回だけ作る。}
\bkword{広義積分}{$\infty$ までの積分。$b$ までの定積分の、$b$ を大きくした先。}
\bkword{ステップ関数 $H(t)$}{$t<0$ で0、$t\geqq0$ で1。時刻0にスイッチを入れる。$\mathcal L[H]=\frac{1}{s}$。}
\end{kotoba}

\begin{column}{ラプラスとヘビサイド：翻訳機ができるまで}
$e^{-st}$ をかけて積分する形は、オイラー（1744年ごろ）やラプラス（1782年ごろから）が調べました。約100年後、イギリスの電気技術者ヘビサイドは、1880年代から1900年ごろにかけて、微分を文字のかけ算として計算する方法（演算子法）で電信や電話の回路を次々に計算しました（6-12 のコラム）。なぜそれでよいのかは説明されず、数学者に疑われましたが、ブロムウィッチ（1916年）たちが数学として整え、ドェッチ（1937年の本）たちが広めました。第二次世界大戦のころから、制御や回路の設計で広く使われています。
\end{column}

\begin{senmonsho}[ラプラス変換の専門書ふうの書き方]
\begin{bkhonbun}
$t\ge0$ で定義された関数 $f(t)$ に対し，
\[ F(s)=\mathcal{L}\{f(t)\}=\int_0^\infty f(t)e^{-st}\,dt \]
をラプラス変換という．$|f(t)|\le Me^{\alpha t}$ ならば，$F(s)$ は $\mathrm{Re}\,s>\alpha$ で存在し，$\mathcal{L}\{f'(t)\}=sF(s)-f(0)$ が成り立つ．
\end{bkhonbun}
\bkyaku{$\mathcal{L}\{\ \}$ は、この本の $\mathcal L[\ ]$（かっこの形がちがうだけ）。「定義された」は「値が決められている」。$\mathrm{Re}\,s$ は $s$ の実部 $\sigma$。「存在し」は「積分が有限の値になり」。$\alpha$ は、この本の「$s>a$」の $a$。$dt$ は、この本の $\dd t$。$f(0)$ を $f(0+)$ や $f(0-)$ と書く本もある（9-6 のせいかくに言うと）。}
\end{senmonsho}




\begin{taisetsu}
\kixfbox{ラプラス変換の3つの段}
\kixfdan{1}\enspace 記号を読む：何に（$f(t)$）・何をかけて（$e^{-st}$）・どこからどこまで（0〜$\infty$）\par
\kixfdan{3}\enspace 直す前に：いくつかの項の和なら、項ごとに（線形）\par
\kixfdan{2}\enspace 表を引く。表の行は、定義の積分で1回だけ作る\par
\[ \kixfNt{\mathcal L[f(t)]}=\kixfSt{1}{\int_0^{\infty}f(t)\,e^{-st}\,\dd t} \]
\kixfbox{場合分けの表9-6（はじめの3行のまとめ）}
\begin{bktable}{colspec={Q[c,m,9mm] Q[l,m,44mm] X[l,m]}, row{2-Z}={font=\small}, rowsep=2pt, colsep=3pt}
記号 & こんなとき & 直し方 \\
\kixfcase{A} & 辞書（表）にある形 & 表を引く（$1\to\frac{1}{s}$、$t\to\frac{1}{s^2}$、$e^{at}\to\frac{1}{s-a}$） \\
\kixfcase{B} & 表の形の和・定数倍 & 項ごとに引いて、たす（線形） \\
\kixfcase{E} & $f'$・$f''$ のある式 & $sF-f(0)$、$s^2F-sf(0)-f'(0)$ \\
\end{bktable}
\end{taisetsu}

%% ---------------- 例題 ----------------
\bksection{reidai}{例題でたしかめよう}

\noindent 例題の「考え方」の1行めには、場合分けの表9-6のどの場合かを書きました（流れ図の道すじは、表\ref{bk:9-5:tab-michi}）。

\begin{reidai}[辞書の行を、定義の積分で作る]
次のラプラス変換を、定義の積分で求めましょう。どんな $s$ のときに使えるかも答えましょう。（1）$\mathcal L[e^{-2t}]$\quad（2）$\mathcal L[3t]$\par
\kangaekata \kixfbaai{表の行を作る}{9-6}{A}\par
\kixfdan{1}\enspace $\mathcal L$：$f$ に $e^{-st}$ をかけて、0から $\infty$ までたす。\par
\kixfnaosu{1}\enspace 0から $b$ までの定積分を計算し、$b$ を大きくする。（2）は部分積分。\par
\shiki\par\nobreak
\begin{kixfhenkei}[0.58]
\int_0^{b}e^{-(s+2)t}\,\dd t=\frac{1}{s+2}\bigl(1-e^{-(s+2)b}\bigr) & （1）肩をまとめて積分（5-14） \\
\to\frac{1}{s+2} & $s+2>0$（$s>-2$）なら $e^{-(s+2)b}\to0$ \\
3\int_0^{\infty}t\,e^{-st}\,\dd t & （2）3を外に出す \\
=3\Bigl[-\frac{t}{s}e^{-st}\Bigr]_0^{\infty}+\frac{3}{s}\int_0^{\infty}e^{-st}\,\dd t & 部分積分（$f=t$、$g'=e^{-st}$） \\
=3\times0+\frac{3}{s}\times\frac{1}{s}=\frac{3}{s^2} & はしの項は0（5-18）。$\mathcal L[1]=\frac{1}{s}$ \\
\end{kixfhenkei}
\kotae （1）\nobreak $\dfrac{1}{s+2}$（$s>-2$）　（2）\nobreak $\dfrac{3}{s^2}$（$s>0$）
\end{reidai}

\begin{jibun}
次のラプラス変換を、定義の積分で求め、どんな $s$ のときに使えるかも答えましょう。（1）$\mathcal L[e^{4t}]$\quad（2）$\mathcal L[4t]$
\bkjans{（1）\nobreak $\int_0^be^{-(s-4)t}\,\dd t=\frac{1}{s-4}(1-e^{-(s-4)b})\to\frac{1}{s-4}$（$s>4$）。（2）\nobreak $4\Bigl[-\frac{t}{s}e^{-st}\Bigr]_0^\infty+\frac{4}{s}\int_0^\infty e^{-st}\,\dd t=\frac{4}{s^2}$（$s>0$）。}
\end{jibun}

\begin{reidai}[冷めるコーヒー]
90\,℃のコーヒーを25\,℃の部屋に置くと、温度 $T$〔℃〕は次の式にしたがいます（$t$ は分）。$T$ をラプラス変換で求め、10分後の温度を求めましょう（$e^{-0.4}=0.6703$）。
\[ \frac{\dd T}{\dd t}=-0.04\,(T-25),\qquad T(0)=90 \]
\kangaekata \kixfbaai{微分方程式を変換する}{9-6}{E}\par
\kixfdan{1}\enspace 両辺の $\mathcal L$ をとる：両辺に $e^{-st}$ をかけて、0から $\infty$ までたす。\par
\kixfnaosu{2}\enspace 微分は $sT(s)-90$、定数1（$0.04\times25$）は $\tfrac{1}{s}$。\par
\shiki\par\nobreak
\begin{kixfhenkei}[0.5]
sT(s)-90=-0.04\,T(s)+\frac{1}{s} & 右辺は $-0.04T+1$ \\
(s+0.04)\,T(s)=90+\frac{1}{s} & $T(s)$ の項を左へ、くくる \\
T(s)=\frac{90}{s+0.04}+\frac{1}{s(s+0.04)} & $s+0.04$ でわる \\
\frac{1}{s(s+0.04)}=\frac{25}{s}-\frac{25}{s+0.04} & 部分分数：$s=0$ で $A=25$、$s=-0.04$ で $B=-25$ \\
T(s)=\frac{25}{s}+\frac{65}{s+0.04} & $90-25=65$ \\
T=25+65e^{-0.04t} & 辞書を逆に引く \\
\end{kixfhenkei}
\kotae $T=25+65e^{-0.04t}$、10分後は $25+65\times0.6703\approx68.6$\,℃
\end{reidai}


\begin{jibun}
70\,℃のお湯を20\,℃の部屋に置きます。温度は $\tfrac{\dd T}{\dd t}=-0.1(T-20)$、$T(0)=70$（$t$ は分）にしたがいます。ラプラス変換で $T$ を求め、5分後の温度を求めましょう（$e^{-0.5}=0.6065$）。
\bkjans{$sT(s)-70=-0.1T(s)+\frac{2}{s}$。$T(s)=\frac{70}{s+0.1}+\frac{2}{s(s+0.1)}$。$\frac{2}{s(s+0.1)}=\frac{20}{s}-\frac{20}{s+0.1}$。$T(s)=\frac{20}{s}+\frac{50}{s+0.1}$、$T=20+50e^{-0.1t}$。5分後は $20+50\times0.6065\approx50.3$\,℃。}
\end{jibun}

\begin{reidai}[RC 回路の充電]
6-4 の回路（電池 \qty{10}{V}、$R=\qty{20}{k\ohm}$、$C=\qty{100}{\micro F}$、時定数 $RC=\qty{2}{s}$）で、コンデンサーの電圧 $V$〔V〕は次の式にしたがいます。ラプラス変換で $V$ を求め、2秒後の電圧を求めましょう（$e^{-1}=0.3679$）。
\[ 2\frac{\dd V}{\dd t}+V=10,\qquad V(0)=0 \]
\kangaekata \kixfbaai{微分方程式を変換する}{9-6}{E}\par
\kixfdan{1}\enspace 両辺の $\mathcal L$ をとる。\par
\kixfnaosu{2}\enspace 微分は $sV(s)-0$、$V$ は $V(s)$、定数10は $\tfrac{10}{s}$。\par
\shiki\par\nobreak
\begin{kixfhenkei}[0.5]
2\bigl(sV(s)-0\bigr)+V(s)=\frac{10}{s} & 1項ずつ辞書で \\
V(s)=\frac{10}{s(2s+1)}=\frac{5}{s(s+0.5)} & $2s+1$ でわり、分母・分子を2でわる \\
\frac{5}{s(s+0.5)}=\frac{10}{s}-\frac{10}{s+0.5} & 部分分数：$s=0$ で $A=10$、$s=-0.5$ で $B=-10$ \\
V=10-10e^{-0.5t} & 辞書を逆に引く \\
V(2)=10-10\times0.3679=6.321 & $t=2$ を入れる \\
\end{kixfhenkei}
\kotae $V=10-10e^{-t/2}$、2秒後は約 \qty{6.32}{V}（6-4 の答えと同じ）
\end{reidai}

\begin{jibun}
電池 \qty{5}{V}・時定数 \qty{1}{s} の RC 回路の電圧は $\tfrac{\dd V}{\dd t}+V=5$、$V(0)=0$ にしたがいます。$V$ を求め、3秒後の電圧を求めましょう（$e^{-3}=0.0498$）。
\bkjans{$sV(s)+V(s)=\frac{5}{s}$。$V(s)=\frac{5}{s(s+1)}=\frac{5}{s}-\frac{5}{s+1}$。$V=5-5e^{-t}$。3秒後は $5-5\times0.0498\approx\qty{4.75}{V}$。}
\end{jibun}

\begin{reidai}[コイルの回路（RL 回路）]
6-4 の RL 回路（電池 \qty{10}{V}、$R=\qty{10}{\ohm}$、$L=\qty{0.5}{H}$）の電流 $i$〔A〕は次の式にしたがいます。$i$ を求め、0.1秒後の電流を求めましょう（$e^{-2}=0.1353$）。
\[ 0.5\,\frac{\dd i}{\dd t}+10\,i=10,\qquad i(0)=0 \]
\kangaekata \kixfbaai{微分方程式を変換する}{9-6}{E}\par
\kixfdan{1}\enspace 両辺の $\mathcal L$ をとる。\par
\kixfnaosu{2}\enspace 微分は $sI(s)-0$、$i$ は $I(s)$、定数10は $\tfrac{10}{s}$。\par
\shiki\par\nobreak
\begin{kixfhenkei}[0.5]
0.5\,sI(s)+10\,I(s)=\frac{10}{s} & 1項ずつ辞書で \\
I(s)=\frac{10}{s(0.5s+10)}=\frac{20}{s(s+20)} & わって、分母・分子を2倍 \\
\frac{20}{s(s+20)}=\frac{1}{s}-\frac{1}{s+20} & 部分分数（$s=0$、$-20$ を代入） \\
i=1-e^{-20t} & 辞書を逆に引く \\
\end{kixfhenkei}
\kotae $i=1-e^{-20t}$、0.1秒後は $1-0.1353\approx\qty{0.865}{A}$
\end{reidai}

\begin{jibun}
RL 回路の電流 $i$〔A〕が $0.2\tfrac{\dd i}{\dd t}+4i=12$、$i(0)=0$ にしたがうとき、$i$ を求め、0.05秒後の電流を求めましょう（$e^{-1}=0.3679$）。
\bkjans{$0.2sI(s)+4I(s)=\frac{12}{s}$。$I(s)=\frac{60}{s(s+20)}=\frac{3}{s}-\frac{3}{s+20}$。$i=3-3e^{-20t}$。0.05秒後は $3-3\times0.3679\approx\qty{1.90}{A}$。}
\end{jibun}

\begin{reidai}[車A を翻訳する]
車A の式とはじめの値は、次のとおりです。ラプラス変換で $X(s)$ の式にし、分母が0になる $s$ を求めましょう。また、部分分数に分けて $x$ を求めましょう。
\[ \ddot x+3\dot x+2x=0,\qquad x(0)=1,\quad \dot x(0)=0 \]
\kangaekata \kixfbaai{微分方程式を変換する}{9-6}{E}\par
\kixfdan{1}\enspace 両辺の $\mathcal L$ をとる。\par
\kixfnaosu{2}\enspace $\ddot x$ は $s^2X-s\cdot1-0$、$\dot x$ は $sX-1$、$x$ は $X$。\par
\shiki\par\nobreak
\begin{kixfhenkei}[0.5]
(s^2X-s)+3(sX-1)+2X=0 & 1項ずつ辞書で \\
(s^2+3s+2)\,X=s+3 & $X$ でくくり、$-s-3$ を右へ \\
X(s)=\frac{s+3}{(s+1)(s+2)} & 分母を因数分解。0になるのは $s=-1$、$-2$ \\
\frac{s+3}{(s+1)(s+2)}=\frac{2}{s+1}-\frac{1}{s+2} & 部分分数：$s=-1$ で $A=2$、$s=-2$ で $B=-1$ \\
x=2e^{-t}-e^{-2t} & 辞書を逆に引く \\
\end{kixfhenkei}
確かめ：$x(0)=2-1=1$、$\dot x(0)=-2+2=0$ で、はじめの値に合っています。\par
はじめの位置が $x(0)=2$ なら、$\ddot x$ は $s^2X-2s-0$、$\dot x$ は $sX-2$ になります（$s\,x(0)=2s$。2をかけ忘れない）。\par
\kotae $X(s)=\dfrac{s+3}{s^2+3s+2}$、分母が0は $s=-1$、$-2$（6-12 の特性方程式の解）、$x=2e^{-t}-e^{-2t}$
\end{reidai}

\begin{jibun}
$\ddot x+5\dot x+6x=0$、はじめ $x(0)=2$、$\dot x(0)=0$ を、ラプラス変換で $X(s)$ の式にし、分母が0になる $s$ を求め、部分分数に分けて $x$ を求めましょう。
\bkjans{$(s^2X-2s)+5(sX-2)+6X=0$。$X(s)=\frac{2s+10}{(s+2)(s+3)}$。分母が0は $s=-2$、$-3$。$\frac{2s+10}{(s+2)(s+3)}=\frac{6}{s+2}-\frac{4}{s+3}$（$s=-2$ で6、$s=-3$ で $\frac{4}{-1}=-4$）。$x=6e^{-2t}-4e^{-3t}$。}
\end{jibun}

\begin{reidai}[微分の行を、数で確かめる]
$f(t)=e^{-5t}$ について、$\mathcal L[f']$ を、（1）$f'$ を先に計算してから辞書で、（2）微分の行 $sF(s)-f(0)$ で求め、同じになることを確かめましょう。\par
\kangaekata \kixfbaai{$f'$ の変換}{9-6}{E}\par
\kixfdan{1}\enspace $\mathcal L[f']$：$f'$ に $e^{-st}$ をかけて、0から $\infty$ までたす。\par
\kixfnaosu{2}\enspace （1）は $e^{at}$ の行、（2）は微分の行を引く。\par
\shiki\par\nobreak
\begin{kixfhenkei}[0.5]
\mathcal L[-5e^{-5t}]=\frac{-5}{s+5} & （1）$f'=-5e^{-5t}$。$e^{at}$ の行の $-5$ 倍 \\
\frac{s}{s+5}-1=\frac{s-(s+5)}{s+5}=\frac{-5}{s+5} & （2）$F=\frac{1}{s+5}$、$f(0)=1$。通分 \\
\end{kixfhenkei}
\kotae どちらも $\dfrac{-5}{s+5}$ で、同じになる
\end{reidai}

\begin{jibun}
$f(t)=e^{3t}$ について、$\mathcal L[f']$ を、（1）$f'$ を先に計算してから辞書で、（2）$sF(s)-f(0)$ で求め、同じになることを確かめましょう。
\bkjans{（1）\nobreak $f'=3e^{3t}$ で $\frac{3}{s-3}$。（2）\nobreak $F=\frac{1}{s-3}$、$f(0)=1$ で $\frac{s}{s-3}-1=\frac{3}{s-3}$。同じ。}
\end{jibun}

\begin{reidai}[たし算の形を変換する]
$\mathcal L[3-2e^{-t}+4t]$ を求めましょう。\par
\kangaekata \kixfbaai{いくつかの項の和}{9-6}{B,A}\par
\kixfdan{1}\enspace $\mathcal L$：$3-2e^{-t}+4t$ に $e^{-st}$ をかけて、0から $\infty$ までたす。\par
\kixfdan{3}\enspace 線形なので、項ごとに辞書を引いてたす。積分はしない。\par
\kixfnaosu{2}\enspace $1\to\tfrac{1}{s}$、$e^{-t}\to\tfrac{1}{s+1}$、$t\to\tfrac{1}{s^2}$。\par
\shiki $3\times\dfrac{1}{s}-2\times\dfrac{1}{s+1}+4\times\dfrac{1}{s^2}$\par
\kotae $\dfrac{3}{s}-\dfrac{2}{s+1}+\dfrac{4}{s^2}$
\end{reidai}

\begin{jibun}
$\mathcal L[5+2t-3e^{2t}]$ を求めましょう。
\bkjans{$\frac{5}{s}+\frac{2}{s^2}-\frac{3}{s-2}$（$s>2$ のとき）。}
\end{jibun}


%% ---------------- まちがえやすいところ ----------------
\bksection{tsumazuki}{まちがえやすいところをチェック}

\begin{tsumazuki}{① 微分の行の $-f(0)$ を忘れる}
\bkng{$\mathcal L\Bigl[\dfrac{\dd T}{\dd t}\Bigr]=sT(s)$}
\bkok{$\mathcal L\Bigl[\dfrac{\dd T}{\dd t}\Bigr]=sT(s)-T(0)$}
$-T(0)$ は、部分積分のはしの項です。忘れると、はじめの値が答えに入らず、お湯の答えが「ずっと20\,℃」のようなまちがいになります。
\end{tsumazuki}

\begin{tsumazuki}{② 記号 $\mathcal L$ を「かける数」と思う}
\bkng{$\mathcal L[3t]=\mathcal L\times3t$}
\bkok{名前を、ふつうの定積分か表に直す}
\[ \kixfNt{\mathcal L[3t]}=\kixfSt{1}{\int_0^{\infty}3t\,e^{-st}\,\dd t}=\frac{3}{s^2} \]
$\mathcal L$ は数ではなく、「$e^{-st}$ をかけて0から $\infty$ までたす」操作の名前です。
\end{tsumazuki}

\begingroup\tcbset{equal height group=kixfEHT}%
\noindent\begin{minipage}[t]{0.485\linewidth}\vspace{0pt}
\begin{tsumazuki}{③ かけ算もそのまま通る}
\raggedright\bkng{$\mathcal L[t\,e^{-t}]=\dfrac{1}{s^2}\times\dfrac{1}{s+1}$}
\bkok{$\mathcal L[fg]$ は $F\times G$ ではない（$\mathcal L[te^{-t}]=\tfrac{1}{(s+1)^2}$：9-6）}
そのまま通るのは、たし算と定数倍だけです。
\end{tsumazuki}
\end{minipage}\hfill
\begin{minipage}[t]{0.485\linewidth}\vspace{0pt}
\begin{tsumazuki}{④ 定数を定数のままにする}
\raggedright\bkng{$\mathcal L[10]=10$}
\bkok{$\mathcal L[10]=\dfrac{10}{s}$}
定数も、0から $\infty$ までたすので、面積は $\tfrac{10}{s}$ です。
\end{tsumazuki}
\end{minipage}\par
\endgroup

\begingroup\tcbset{equal height group=kixfEHU}%
\noindent\begin{minipage}[t]{0.485\linewidth}\vspace{0pt}
\begin{tsumazuki}{⑤ $s$ の条件を考えない}
\raggedright\bkng{$\mathcal L[e^{4t}]$ で $s=2$ にする}
\bkok{使えるのは $s>4$ だけ}
\end{tsumazuki}
\end{minipage}\hfill
\begin{minipage}[t]{0.485\linewidth}\vspace{0pt}
\begin{tsumazuki}{⑥ 小文字と大文字をまぜる}
\raggedright\bkng{${sx-x(0)+x=\dfrac{1}{s}}$}
\bkok{${sX(s)-x(0)+X(s)=\dfrac{1}{s}}$}
\end{tsumazuki}
\end{minipage}\par
\endgroup

%% ---------------- 物理・工学ではこう使う ----------------
\begin{butsuri}
\kixfbox{スイッチを入れた直後の電気回路（過渡現象）}
\noindent\begin{minipage}[c]{0.6\linewidth}
\raggedright
スイッチを入れたり切ったりした直後の、電圧や電流の変わり方を\textbf{過渡現象}（かとげんしょう）といいます。抵抗・コイル・コンデンサーの回路の式は、$L\tfrac{\dd i}{\dd t}$、$\tfrac{1}{C}\int i\,\dd t$ をふくむ微分方程式です。ラプラス変換すると、コイルは「$sL$」、コンデンサーは「$\tfrac{1}{sC}$」（はじめの電流と電荷が0のとき）という、抵抗と同じように計算できる量になり、回路をオームの法則のように解けます（物理 電-9）。
\end{minipage}\hfill
\begin{minipage}[c]{0.37\linewidth}
\centering
\begin{tikzpicture}[x=0.85mm, y=1mm]
  \draw[bkInk, line width=0.8pt] (0,10) -- (0,18) -- (8,18);
  \draw[bkInk, line width=0.8pt] (0,7.5) -- (0,0);
  \draw[bkInk, line width=0.8pt] (8,18) -- (15,21);
  \fill[bkInk] (8,18) circle (0.7) (16,18) circle (0.7);
  \draw[bkInk, line width=0.8pt] (16,18) -- (20,18);
  \draw[bkInk, line width=0.8pt, fill=white] (20,16) rectangle (30,20);
  \draw[bkInk, line width=0.8pt] (30,18) -- (36,18) -- (36,11);
  \draw[bkInk, line width=1.1pt] (32,11) -- (40,11) (32,8) -- (40,8);
  \draw[bkInk, line width=0.8pt] (36,8) -- (36,0) -- (0,0);
  \draw[bkInk, line width=1.1pt] (-3,10) -- (3,10);
  \draw[bkInk, line width=0.7pt] (-1.6,7.5) -- (1.6,7.5);
  \node[font=\footnotesize, anchor=south] at (25,21) {$R$};
  \node[font=\footnotesize, anchor=west] at (41,9.5) {$C$};
  \node[font=\footnotesize, anchor=south] at (12,22) {スイッチ};
\end{tikzpicture}
\end{minipage}
\par\vspace{4pt}
\kixfbox{制御の設計}
エアコン・エレベーター・ドローンなど、目標に合わせて機械を動かす「制御」では、機械の式をラプラス変換して $s$ の式で表し、かけ算・わり算で設計します（9-8・9-9、物理 波-4）。微分方程式を毎回解かなくてよいので、たくさんの部品をつないだ大きな機械でも計算できます。
\par\vspace{4pt}
\kixfbox{電信と電話の回路}
ヘビサイドが演算子法を使ったのは、遠くまで信号を送る電信・電話の線（送電線と同じ、分布定数回路）の計算でした。信号がゆがまずに届く条件（ヘビサイドの条件）も見つけています。
\end{butsuri}

%% ---------------- 章末問題 ----------------
\bkendproblems
\bkproblemgroup{基本}{5}

\bkproblem $\mathcal L[e^{-3t}]$ を、定義の積分で求めましょう。どんな $s$ のときに使えますか\mbox{（\bkreidairef{1}）}。
\bkpans{$\int_0^be^{-(s+3)t}\,\dd t\to\frac{1}{s+3}$。$s>-3$ のとき。}

\bkproblem $\mathcal L[5t]$ を、定義の積分と部分積分で求めましょう\mbox{（\bkreidairef{1}）}。
\bkpans{$5\Bigl[-\frac{t}{s}e^{-st}\Bigr]_0^\infty+\frac{5}{s}\int_0^\infty e^{-st}\,\dd t=\frac{5}{s^2}$（$s>0$）。}

\bkproblem $f(t)=e^{-3t}$ について、$\mathcal L[f']$ を、$f'$ を先に計算する方法と、$sF(s)-f(0)$ の方法で求め、同じになることを確かめましょう\mbox{（\bkreidairef{6}）}。
\bkpans{$f'=-3e^{-3t}$ で $\frac{-3}{s+3}$。$\frac{s}{s+3}-1=\frac{-3}{s+3}$。同じ。}

\bkproblem $\mathcal L[3+2t-e^{-4t}]$ を求めましょう\mbox{（\bkreidairef{7}）}。
\bkpans{$\frac{3}{s}+\frac{2}{s^2}-\frac{1}{s+4}$。}

\bkproblem 60\,℃のお湯を20\,℃の部屋に置きます。温度は $\tfrac{\dd T}{\dd t}=-0.05(T-20)$（$t$ は分）にしたがいます。ラプラス変換で $T$ を求め、10分後の温度を求めましょう（$e^{-0.5}=0.6065$）\mbox{（\bkreidairef{2}）}。
\bkpans{$sT(s)-60=-0.05T(s)+\frac{1}{s}$。$T(s)=\frac{60}{s+0.05}+\frac{20}{s}-\frac{20}{s+0.05}=\frac{20}{s}+\frac{40}{s+0.05}$。$T=20+40e^{-0.05t}$。10分後は約44.3\,℃。}

\bkproblem $f(t)=e^{2t}$ について、$\mathcal L[f']$ を、$f'$ を先に計算する方法と、$sF(s)-f(0)$ の方法で求め、同じになることを確かめましょう\mbox{（\bkreidairef{6}）}。
\bkpans{$f'=2e^{2t}$ で $\frac{2}{s-2}$。$\frac{s}{s-2}-1=\frac{2}{s-2}$。同じ。}

\bkproblem $\mathcal L[6-t+2e^{3t}]$ を求めましょう\mbox{（\bkreidairef{7}）}。
\bkpans{$\frac{6}{s}-\frac{1}{s^2}+\frac{2}{s-3}$（$s>3$ のとき）。}

\bkproblemgroup{チャレンジ}{5}

\bkproblem 電池 \qty{12}{V}、$R=\qty{10}{k\ohm}$、$C=\qty{100}{\micro F}$ の回路で、コンデンサーの電圧は $\tfrac{\dd V}{\dd t}+V=12$、$V(0)=0$ です（時定数 \qty{1}{s}）。ラプラス変換で $V$ を求め、2秒後の電圧を求めましょう（$e^{-2}=0.1353$）\mbox{（\bkreidairef{3}）}。
\bkpans{$V(s)=\frac{12}{s(s+1)}=\frac{12}{s}-\frac{12}{s+1}$。$V=12-12e^{-t}$。2秒後は約 \qty{10.4}{V}。}

\bkproblem 電池 \qty{6}{V}、$R=\qty{3}{\ohm}$、$L=\qty{0.3}{H}$ の RL 回路で、電流 $i$〔A〕は $0.3\tfrac{\dd i}{\dd t}+3i=6$、$i(0)=0$ にしたがいます。ラプラス変換で $i$ を求め、0.1秒後の電流を求めましょう（$e^{-1}=0.3679$）\mbox{（\bkreidairef{4}）}。
\bkpans{$0.3sI(s)+3I(s)=\frac{6}{s}$。$I(s)=\frac{20}{s(s+10)}=\frac{2}{s}-\frac{2}{s+10}$。$i=2-2e^{-10t}$。0.1秒後は $2-2\times0.3679\approx\qty{1.26}{A}$。}

\bkproblem $\ddot x+4\dot x+3x=0$、はじめ $x(0)=2$、$\dot x(0)=0$ を、ラプラス変換で $X(s)$ の式にし、分母が0になる $s$ を求め、部分分数に分けて $x$ を求めましょう\mbox{（\bkreidairef{5}）}。
\bkpans{$(s^2X-2s)+4(sX-2)+3X=0$。$X(s)=\frac{2s+8}{(s+1)(s+3)}$。分母が0は $s=-1$、$-3$。$\frac{3}{s+1}-\frac{1}{s+3}$（$s=-1$ で $\frac{6}{2}=3$、$s=-3$ で $\frac{2}{-2}=-1$）。$x=3e^{-t}-e^{-3t}$。}

\begin{furikaeri}
\begin{checklist}
\item 「微分→$s$ をかける」の辞書を使って、冷めるお湯の微分方程式を $s$ の式に書きかえ、表でもどして答えを出せる。\bkfurref{bk:9-5:sekai}
\item ラプラス変換の記号を「名前」として読み、辞書（変換表）の行を、定義の積分で1回だけ作れる。\bkfurref{bk:9-5:gyou}
\item 微分の行 $sF(s)-f(0)$ を部分積分で導き、場合分けの表9-6（A・B・E）を使って微分方程式を変換できる。\bkfurref{bk:9-5:bibun}
\end{checklist}
\end{furikaeri}

\bkfukidashi{sora}{微分が $s$ のかけ算になって、はじめの値も自動で入る。次は、辞書の行をふやして、表を完成させよう。}
